\documentclass[11pt]{article}

\usepackage[a4paper,margin=1in]{geometry}
\usepackage[T1]{fontenc}
\usepackage[utf8]{inputenc}
\usepackage{lmodern}
\usepackage{microtype}
\usepackage{amsmath,amssymb,amsthm,mathtools}
\usepackage{booktabs}
\usepackage{array}
\usepackage{xurl}
\usepackage{hyperref}
\usepackage{authblk}

\allowdisplaybreaks
\numberwithin{equation}{section}

\hypersetup{
  unicode=true,
  colorlinks=true,
  linkcolor=blue,
  citecolor=blue,
  urlcolor=blue,
  pdftitle={A Proposed Complete Solution to Erdos Problem 1038},
  pdfauthor={Shouqiao Wang}
}
\urlstyle{same}

\theoremstyle{plain}
\newtheorem{theorem}{Theorem}[section]
\newtheorem{lemma}[theorem]{Lemma}
\newtheorem{proposition}[theorem]{Proposition}
\newtheorem{corollary}[theorem]{Corollary}
\theoremstyle{definition}
\newtheorem{definition}[theorem]{Definition}
\newtheorem{certificate}[theorem]{Certificate}
\theoremstyle{remark}
\newtheorem{remark}[theorem]{Remark}

\newcommand{\R}{\mathbb R}
\newcommand{\Z}{\mathbb Z}
\newcommand{\T}{\mathbb T}
\newcommand{\1}{\mathbf 1}
\newcommand{\dd}{\,d}
\newcommand{\pv}{\operatorname {pv}}
\newcommand{\supp}{\operatorname {supp}}
\newcommand{\eps}{\varepsilon}
\newcommand{\cA}{\mathcal A}
\newcommand{\J}{\mathfrak J}
\newcommand{\E}{\mathcal E}
\newcommand{\norm}[1]{\left\lVert#1\right\rVert}

\title{
A Proposed Complete Solution to Erd\H{o}s Problem 1038
}
\author{Shouqiao Wang}
\affil{Columbia University}
\date{}

\begin{document}
\maketitle

\begin{abstract}
Let $f$ range over the nonconstant monic real polynomials whose zeros all
belong to $[-1,1]$.  We determine the infimum and supremum of the Lebesgue
measure of $\{x\in\R:|f(x)|<1\}$.  The supremum is $2\sqrt2$.  The
infimum is a uniquely determined constant $L$, characterized below as the
minimum of a one-variable function, with numerical value
\[
 L=1.834430475762661\ldots.
\]
The lower-bound proof combines component atomization, an elementary
mean-deficit estimate, a convex constant-platform comparison, an
endpoint-corrected adjoint identity, and a circle rearrangement theorem.
All remaining scalar and parameter-uniform signs are certified by directed
outward interval arithmetic in one companion Python file.  An explicit
positive-platform approximation proves sharpness.  Every finite polynomial
satisfies the lower inequality strictly, so the infimum is not attained;
the supremum is attained by $(x^2-1)^m$ for every $m\ge1$.
\end{abstract}

\section{Introduction and statement of the result}

Erd\H{o}s, Herzog, and Piranian asked for the two extremal values of
\[
 \left|\{x\in\R:|f(x)|<1\}\right|
\]
when $f$ is a nonconstant monic polynomial with all of its zeros real and
contained in $[-1,1]$~\cite[p.~131]{EHP1958}; see also the modern
formulation~\cite{Bloom1038}.  We give a computer-assisted determination of both
extrema.  The proof of the supremum follows the expansive-rearrangement
argument in Tao's note~\cite{Tao1038}; we include the full argument and
independently certify its three scalar inequalities.  The lower bound and
the sharp limiting construction are proved in
Sections~\ref{sec:reduction}--\ref{sec:recovery} below.

The endpoint-atom reductions and the one-cut limiting candidate were also
developed in the public discussion surrounding the
problem~\cite{Bloom1038Discussion}.  The new burden addressed here is the global
lower comparison for every finite root configuration, together with one
exhaustive certificate for the remaining parameter inequalities.  Thus we
do not present the supremum argument or the initial numerical candidate as
new; the claimed result is the complete lower-bound closure and its
reproducible verification.

Let $\mathcal P$ be the set of all nonconstant monic real polynomials whose
zeros, counted with multiplicity, lie in $[-1,1]$, and put
\[
 E_f=\{x\in\R:|f(x)|<1\}.
\]
For $0<q<3-2\sqrt2$ define
\[
 H(q)=\frac{2q}{(1+q)^2},\qquad
 s(q)=\frac{1-q}{1+q},\qquad
 A(q)=\frac{\log H(q)}{\log q}.
\tag{1.1}
\]
Let $q_s$ be the unique solution in $(0,3-2\sqrt2)$ of
$A(q_s)=s(q_s)$.  For
$0<q\le q_s$, let $u_-(q)$ be the unique solution of
\[
 A(q)\log\frac{u-q}{|1-qu|}=\log u
\tag{1.2}
\]
in $u>q^{-1}$.  For $0<q<q_s$, let $u_+(q)$ be the unique nontrivial
solution of (1.2) in $1<u<q^{-1}$.  At $q=q_s$ set $u_+(q_s)=1$ by
continuity.  Define
\[
 \Lambda(q)=H(q)\left(u_-(q)+u_-(q)^{-1}
                         -u_+(q)-u_+(q)^{-1}\right).
\tag{1.3}
\]

\begin{theorem}[Main theorem]\label{thm:main}
The function $\Lambda$ has a unique minimizer $q_*$ on $(0,q_s]$.  If
$L=\Lambda(q_*)$, then
\[
 \boxed{\inf_{f\in\mathcal P}|E_f|=L,
        \qquad \sup_{f\in\mathcal P}|E_f|=2\sqrt2.}
\tag{1.4}
\]
The infimum is not attained.  The supremum is attained by
$(x^2-1)^m$ for every integer $m\ge1$.

The following outward enclosures hold:
\[
\begin{aligned}
 0.025715536866527&<q_*<0.025715536866528,\\
 1.834430475762661&<L<1.834430475762662.
\end{aligned}
\]
\end{theorem}

\paragraph{The exact constant and its computation.}
The symbol $L$ denotes the single exact real number $\Lambda(q_*)$, not
the numerical interval displayed above.  That interval has deliberately
been rounded outward and shortened for readability; it is only a rigorous
enclosure of $L$.  The value is computed by certified one-dimensional root
isolation, not by minimizing over a sampled grid.  Indeed, put
\[
 F(q,u)=A(q)\log\frac{u-q}{|1-qu|}-\log u,
 \qquad W(u)=u+u^{-1}.
\]
For $0<q<q_s$, the two roots $u_+(q)$ and $u_-(q)$ lie on the unique
branches specified above.  On either branch the implicit-function theorem
gives
\[
 u_\pm'(q)=-\frac{\partial_qF(q,u_\pm(q))}
                  {\partial_uF(q,u_\pm(q))},
\]
and hence
\[
 \Lambda'(q)
 =H'(q)\bigl(W(u_-)-W(u_+)\bigr)
 +H(q)\bigl(W'(u_-)u_-'-W'(u_+)u_+'\bigr).
\tag{1.5}
\]
First one isolates $q_s$, equivalently as the unique zero of the strictly
decreasing function in (6.4).  For $0<q<q_s$, one isolates the two
nondegenerate roots of $F(q,u)=0$; at $q=q_s$, one has $u_+=1$ exactly,
and the regularized variable $y=(\log u_+)^2$ in (6.12)--(6.12b) gives
the continuous soft-edge branch.  One then isolates the unique root
$q_*$ of $\Lambda'(q)=0$ and evaluates $\Lambda(q_*)$.  Certified
bisection or interval Newton iteration, applied to the nondegenerate or
regularized equations as appropriate, makes all of these enclosing
intervals arbitrarily narrow.  The companion verifier certifies the
particular outward enclosures displayed here.

The limiting probability measure producing the lower constant is explicit.
Put $A_*=A(q_*)$, $s_*=s(q_*)$, and $a_*=2s_*^2-1$.  Then
\[
 \mu_*=A_*\delta_{-1}+
 \frac{t+1-2A_*s_*}
 {\pi(t+1)\sqrt{(t-a_*)(1-t)}}\1_{[a_*,1]}(t)\dd t.
\tag{1.6}
\]
Its negative-potential component has endpoints
\[
 \alpha_*=s_*^2-H(q_*)\bigl(u_-(q_*)+u_-(q_*)^{-1}\bigr),\qquad
 \beta_*=s_*^2-H(q_*)\bigl(u_+(q_*)+u_+(q_*)^{-1}\bigr),
\tag{1.7}
\]
and $\beta_*-\alpha_*=L$.

We briefly describe the proof and its trust boundary.  A polynomial is first
replaced by its empirical root measure.  Simultaneous barycentric collapse in
the connected components of the sublevel set cannot increase its length and
reduces the lower-bound problem to a finite configuration consisting of one
endpoint atom and finitely many residual atoms.  A direct energy estimate
handles the small endpoint-to-residual mass ratio.  For all larger ratios, a
constant-potential reference measure gives a supporting inequality for the
main component; a circle rearrangement theorem then reduces every residual
block to explicit one-variable inequalities.  Finally, a zero-platform
one-cut family identifies $L$, and positive-platform empirical measures
approach it from above.

The analytic reductions, including all limiting and strictness arguments,
are proved in the paper.  Three explicitly stated certificates remain: the
global one-cut minimization, the uniform parameter cover, and the scalar
signs in the upper-bound argument.  The single companion file
\texttt{numerical\_verifier.py} checks precisely those claims with outward
rounding; it does not search for a proof and it aborts whenever a required
sign cannot be resolved.  Appendix~\ref{app:verification} records the exact
scope of this computation.

\section{Potentials, atomization, and normalization}\label{sec:reduction}

For a probability measure $\mu$ on $[-1,1]$ write
\[
 V_\mu(x)=\int\log|x-t|\dd\mu(t),\qquad
 E_\mu=\{x:V_\mu(x)<0\}.
\tag{2.1}
\]
If $f(x)=\prod_{j=1}^n(x-r_j)$ and
$\mu_f=n^{-1}\sum_j\delta_{r_j}$, then
$V_{\mu_f}=n^{-1}\log|f|$ and $E_f=E_{\mu_f}$.  Moreover
$E_\mu\subset(-2,2)$, because $|x-t|\ge1$ for $|x|\ge2$ and
$t\in[-1,1]$.  We call a measure of the form $\mu_f$ an
\emph{empirical measure}; its atom masses are positive rational numbers
whose sum is one.

\begin{lemma}[Simultaneous component atomization]\label{lem:atomize}
Let $\mu$ be an empirical root measure and write the finite union of
components of $E_\mu$ as $\bigsqcup_j I_j$.  In every $I_j$, replace all
root mass by an atom of the same mass at its barycenter.  If $\widetilde
\mu$ is the resulting empirical measure, then
\[
 E_{\widetilde\mu}\subset E_\mu.
\tag{2.2}
\]
Every component of $E_{\widetilde\mu}$ contains exactly one distinct root
cluster.
\end{lemma}

\begin{proof}
Write the distinct roots as $c_h$, with masses $p_h>0$.  If the degree of
the underlying polynomial is $n$, then $V_\mu=n^{-1}\log|f|$.  Every finite
boundary point of $E_\mu$ is therefore a real zero of the nonzero polynomial
$f^2-1$.  There are only finitely many such points; together with
$E_\mu\subset(-2,2)$, this proves that $E_\mu$ has finitely many components.
Each component $I_j=(a_j,b_j)$ contains a root.  Otherwise
\[
 V_\mu''(x)=-\sum_h\frac{p_h}{(x-c_h)^2}<0
 \qquad (x\in I_j).
\]
The potential is continuous at $a_j,b_j$ and equals zero there.  Strict
concavity would then give
$V_\mu(x)\ge0$ for $a_j<x<b_j$, contrary to the definition of $I_j$.

Fix $x\notin E_\mu$.  Then $x\notin I_j$ for every $j$, and
\[
 \frac{\partial^2}{\partial t^2}\log|x-t|
 =-\frac1{(x-t)^2}<0\qquad(t\in I_j).
\]
For each $j$, Jensen's inequality therefore shows that replacing the root
mass in $I_j$ by its barycenter raises the contribution of that mass to the
potential at $x$.  All roots belong to $E_\mu$ because the potential is
$-\infty$ at a root, so these replacements account for the entire measure.
Doing them simultaneously gives
$V_{\widetilde\mu}(x)\ge V_\mu(x)\ge0$, proving (2.2).

Inside an old component $I_j$, the new potential has only one pole $c_j$.
On each of $(a_j,c_j)$ and $(c_j,b_j)$ it is strictly concave, tends to
$-\infty$ at $c_j$, and is nonnegative at the other endpoint.  On, say,
$(a_j,c_j)$, a concave function has convex superlevel sets; equivalently,
once it becomes negative while moving from $a_j$ toward $c_j$, it cannot
become nonnegative again.  Its negative set on that side is therefore the
interval adjacent to $c_j$; the same holds on the right.  The two pieces
form one component containing $c_j$.  By (2.2), the complementary gaps
between distinct old components remain nonnegative, so different new
components cannot merge.  Thus every new component contains exactly one
root cluster.
\end{proof}

\subsection{Orientation and endpoint normalization}

We henceforth replace $\mu_f$ by the atomized measure from
Lemma~\ref{lem:atomize}.  This can only decrease the sublevel set, so a
lower bound for the atomized measure is also a lower bound for the original
polynomial.  If $\mu^{-}$ denotes the pushforward of $\mu$ under
$t\mapsto-t$, then $V_{\mu^-}(x)=V_\mu(-x)$ and hence
$|E_{\mu^-}|=|E_\mu|$.  We choose between $\mu$ and $\mu^-$ so that the
root mean
\[
 m=\int t\dd\mu(t)
\]
is nonpositive.  For
$-1<x<0$ and $-1\le t\le1$,
\[
 (1-xt)^2-(x-t)^2=(1-x^2)(1-t^2)\ge0,
\]
so $|x-t|\le1-xt$.  Jensen gives
\[
 V_\mu(x)\le\int\log(1-xt)\dd\mu(t)
 \le\log(1-xm)\le0.
\tag{2.3}
\]
The inequality is strict.  Indeed, equality in Jensen for the strictly
concave function $u\mapsto\log u$ forces $1-xt$ to be constant
$\mu$-almost everywhere, hence $\mu$ to be a point mass.  Equality in the
last inequality forces $m=0$, so that point mass would be $\delta_0$; for
$\delta_0$ the first inequality is strict whenever $-1<x<0$.
Thus
\[
 (-1,0)\subset E_\mu.
\tag{2.4}
\]

Let $c$ be the unique root cluster in the component containing
$(-1,0)$, and let its mass be $A$.  It is the leftmost cluster: any root
to its left would lie in the same component, contradicting the one-cluster
conclusion of Lemma~\ref{lem:atomize}.  Thus every root is at least $c$.
Since their mean is nonpositive, $c\le m\le0$.

Translate every root, and the real variable, by $-c-1$; this sends $c$ to
$-1$ and preserves all lengths.  A root $t\in[c,1]$ is sent to
$t-c-1\in[-1,-c]\subset[-1,1]$.  The new mean is
$m-c-1\le -c-1\le0$, because $m\le0$ and $c\ge-1$.  Denote the translated
measure temporarily by $\mu^{\rm tr}$.

Let $r_0$ be the right endpoint of the main component before this
translation.  By (2.4),
$r_0\ge0$.  If $r_0-c<1$, the distance from $r_0$ to the main cluster is
strictly less than one.  Every residual cluster lies to the right of that
component, hence in $[r_0,1]$, and its distance from $r_0$ is at most one.
The weighted geometric mean of these distances would therefore be strictly
less than one, contradicting $\exp V_\mu(r_0)=1$ for the unshifted measure.
Consequently
\[
 \beta=r_0-c-1\ge0.
\tag{2.5}
\]
We now rename $\mu^{\rm tr}$ as $\mu$.  After translation, the main atom is
at $-1$, the translated right boundary
is $\beta$, and all residual roots belong to $[\beta,1]$; in particular
$0\le\beta\le1$ when residual mass is present.  At this boundary,
weighted AM--GM gives
\[
 1=e^{V_\mu(\beta)}
 \le\int|\beta-t|\dd\mu(t)
 \le A(1+\beta)+(1-A)(1-\beta)
 =1+(2A-1)\beta.
\tag{2.6}
\]
If $\beta>0$, this proves $A\ge1/2$.  If $\beta=0$, equality in the
geometric--arithmetic mean chain forces every distance from $\beta$ to a
residual root to equal one, hence every residual root to be $1$.  The
nonpositive-mean condition then gives
$-A+(1-A)\le0$, again $A\ge1/2$.

If $A=1$, the sublevel interval is $(-2,0)$ and has length $2>L$.
Assume henceforth $B=1-A>0$, and write
\[
 k=\frac AB\ge1,\qquad
 \mu=A\delta_{-1}+B\sum_{i=1}^Nq_i\delta_{t_i},
 \quad 0\le t_1<\cdots<t_N\le1,
 \quad\sum_iq_i=1.
\tag{2.7}
\]
In the shifted spatial coordinate $x_{\rm new}=x_{\rm old}+1$, put
$d_i=1+t_i\in[1,2]$.  Addition of this translation and division of the
potential by the positive number $B$ do not change its sign, and give
\[
 W(x)=k\log|x|+\sum_iq_i\log|x-d_i|.
\tag{2.8}
\]
Let $M_k$ be the width of the component of $\{W<0\}$ containing zero and
define
\[
 R_i=\exp\left[-\frac1{q_i}\left(k\log d_i+
              \sum_{j\ne i}q_j\log|d_i-d_j|\right)\right].
\tag{2.9}
\]

\begin{lemma}[Local component radius]\label{lem:radius}
The component containing $d_i$ has width at least $2R_i$.  Hence
\[
 |E_\mu|\ge \J_k:=M_k+2\sum_iR_i.
\tag{2.10}
\]
\end{lemma}

\begin{proof}
Let $(\ell_i,r_i)$ be that component, set
$s=d_i-\ell_i$, $u=r_i-d_i$, $T=s+u$, and $w=s/T$.  After removing the
self-pole, the background
\[
 H_i(x)=k\log x+\sum_{j\ne i}q_j\log|x-d_j|
\]
is well defined and strictly concave on $(\ell_i,r_i)$.  To justify the
first assertion, for $0<x<1$ we have directly
$W(x)\le k\log x+\log(2-x)<0$.  Thus the main component contains
$(0,1)$, while atomization puts the residual cluster $d_i$ in a different
component; hence $\ell_i\ge1>0$.  The second assertion follows
by differentiating twice after the self-pole has been removed.

The boundary equations and (2.9) are
\[
 H_i(\ell_i)=-q_i\log s,\quad
 H_i(r_i)=-q_i\log u,\quad
 H_i(d_i)=-q_i\log R_i.
\]
Since $d_i=(1-w)\ell_i+wr_i$, concavity implies
\[
 -q_i\log R_i
 \ge-(1-w)q_i\log s-wq_i\log u.
\]
Division by $-q_i<0$ reverses the inequality, and exponentiation gives
\[
 R_i\le s^{1-w}u^w
 =T w^{1-w}(1-w)^w\le T/2.
\]
For completeness, if
$h(w)=(1-w)\log w+w\log(1-w)$, then
$h'(w)=\log((1-w)/w)+(1-2w)/(w(1-w))$ is positive on
$(0,1/2)$ and changes sign by symmetry at $1/2$; hence
$w^{1-w}(1-w)^w\le1/2$.  Thus $T\ge2R_i$.  The main component and all
residual components are pairwise disjoint, so adding their lengths proves
(2.10).
\end{proof}

\begin{lemma}[Endpoint window]\label{lem:window}
Let $\rho(k)\in(0,1)$ be the unique solution of
\[
 \rho(k)^k(2+\rho(k))=1.
\tag{2.11}
\]
Then $M_k\ge1+\rho(k)$, and $\rho(k)\ge\sqrt2-1$ for $k\ge1$.
\end{lemma}

\begin{proof}
For fixed $k\ge1$, the function $r\mapsto r^k(2+r)$ is continuous and
strictly increasing on $[0,1]$, with values $0$ and $3$ at the endpoints.
This proves the existence and uniqueness of $\rho(k)$.
For $0<x<1$, since $d_i-x\le2-x$,
\[
 W(x)\le k\log x+\log(2-x)\le\log(x(2-x))<0.
\]
For $0<r<\rho(k)$,
\[
 W(-r)\le k\log r+\log(2+r)<0.
\]
Thus $(-\rho(k),1)$ lies in the main component.  At $k=1$,
$\rho(1)=\sqrt2-1$; for fixed $0<\rho<1$, the left side of (2.11)
decreases as $k$ increases, so its solution increases.
\end{proof}

\section{The elementary range: a mean-deficit energy estimate}

The following argument is important because it does not use a
fixed-endpoint or unrestricted Bojanov comparison.

\begin{proposition}[Small endpoint-to-residual ratio]\label{prop:lowk}
For every atomized configuration (2.8) with
$1\le k\le K:=29/20$,
\[
 \sum_iR_i>\frac13,\qquad \J_k>2.
\tag{3.1}
\]
\end{proposition}

\begin{proof}
Separation of the main component from the component of $d_1$ supplies
$b\in(0,d_1)$ such that $W(b)\ge0$.  Put
$y_i=2-d_i\ge0$ and $\bar y=\sum_iq_i y_i$.  Jensen's inequality gives
\[
 0\le k\log b+\sum_iq_i\log(d_i-b)
 \le k\log b+\log(2-\bar y-b).
\tag{3.2}
\]
In particular $b^k(2-b)\ge1$.  If $b<1$, then
$b^k(2-b)\le b(2-b)<1$, a contradiction.  Thus $b\ge1$ and
\[
 \bar y\le2-b-b^{-k}
 \le\eps_K:=2-(K+1)K^{-K/(K+1)}.
\tag{3.3}
\]
Indeed, $b\ge1$ and $k\le K$ imply $b^{-k}\ge b^{-K}$, so
$2-b-b^{-k}\le2-b-b^{-K}$.  The derivative of the latter function is
$-1+Kb^{-K-1}$; it vanishes only at $b=K^{1/(K+1)}$, is positive before
that point, and is negative after it.  This proves the displayed maximum.
We now prove the exact rational bound
\[
 0<\eps_K<\frac1{25}.
\tag{3.4}
\]
Positivity follows because the derivative at $b=1$ is $K-1>0$.
Furthermore
\[
 e^{3/8}>1+\frac38+\frac{(3/8)^2}{2}
               +\frac{(3/8)^3}{6}
 =\frac{1489}{1024}>\frac{29}{20},
\]
so $\log(29/20)<3/8$.  With $z=87/392$, Taylor's theorem gives
$e^{-z}>1-z+z^2/2-z^3/6$, and direct rational subtraction gives
\[
 \frac{49}{20}\left(1-z+\frac{z^2}{2}-\frac{z^3}{6}\right)
 -\frac{49}{25}=\frac{522631}{245862400}>0.
\]
Here $K/(K+1)=29/49$ and
$z=(29/49)(3/8)$.  Hence

\[
 K^{-K/(K+1)}
 =\exp\left(-\frac{29}{49}\log K\right)>e^{-z},
\]
so the last rational inequality says
$(K+1)K^{-K/(K+1)}>49/25$.  Therefore
$\eps_K=2-(K+1)K^{-K/(K+1)}<1/25$, proving (3.4).

Put $Q=\sum_iq_i^2$ and $t=1-Q=2\sum_{i<j}q_iq_j$.  For $t>0$,
concavity of the logarithm and $|d_i-d_j|=|y_i-y_j|\le y_i+y_j$ yield
\begin{align*}
2\sum_{i<j}q_iq_j\log|d_i-d_j|
&\le t\log\left(\frac{2\sum_{i<j}q_iq_j|d_i-d_j|}{t}\right)\\
&\le t\log\left(\frac{2\sum_iq_i(1-q_i)y_i}{t}\right)
 \le t\log\frac{2\eps_K}{t}.
\tag{3.5}
\end{align*}
For $t=0$ the pair sums are empty and the terms involving
$t\log(a/t)$ are defined by continuity as zero.

The definition (2.9) gives the exact identity
\[
 \sum_iq_i^2\log R_i
 =-k\sum_iq_i\log d_i
  -2\sum_{i<j}q_iq_j\log|d_i-d_j|.
\tag{3.6}
\]
Because $d_i\le2$, (3.5)--(3.6), followed by weighted AM--GM with
weights $q_i^2/Q$, gives, with $S=\sum_iR_i$,
\[
 S\ge\sum_i\frac{q_i^2}{Q}R_i
 \ge\exp\left(\frac{-k\log2-t\log(2\eps_K/t)}{1-t}\right).
\]
Equivalently,
\[
 (1-t)\log(3S)
 \ge\log3-k\log2-t\log\frac{6\eps_K}{t}.
\tag{3.7}
\]
The elementary maximum $t\log(a/t)\le a/e$, (3.4), and
$e>1+1+1/2+1/6=8/3$ show
\[
 t\log\frac{6\eps_K}{t}<\frac9{100}.
\tag{3.8}
\]
Finally, the positive atanh series gives
\[
 \log3>2\left(\frac12+\frac{(1/2)^3}{3}+\frac{(1/2)^5}{5}\right)
 =\frac{263}{240},
\]
whereas the first two terms and a geometric tail give
\[
 \log2<2\left(\frac13+\frac{(1/3)^3}{3}\right)
 +\frac{2(1/3)^5}{5(1-1/9)}=\frac{1123}{1620}.
\]
Hence
\[
 \log3-\frac{29}{20}\log2
 >\frac{1469}{16200}=\frac9{100}+\frac{11}{16200}.
\tag{3.9}
\]
Equations (3.7)--(3.9) imply
$(1-t)\log(3S)>11/16200$.  Since $1-t=Q=\sum_iq_i^2>0$, this gives
$S>1/3$.  By
Lemma~\ref{lem:window},
\[
 \J_k>\sqrt2+\frac23>2,
\]
because $\sqrt2>4/3$.
\end{proof}

\section{Constant-platform references and convex calibration}

Fix $k\ge1$ and $1\le a<2$, put $I=[a,2]$, and define
\[
 c=\frac{a+2}{2},\qquad r=\frac{2-a}{2},\qquad H=\frac r2=\frac{2-a}{4}.
\tag{4.1}
\]
The equilibrium probability of $I$ and the balayage of $\delta_0$ onto
$I$ are
\[
 \dd e_I(d)=\frac{\dd d}{\pi\sqrt{(d-a)(2-d)}},\qquad
 \dd\omega_{0,I}(d)=\frac{\sqrt{2a}}d\,\dd e_I(d).
\tag{4.2}
\]
These are the standard interval equilibrium and balayage
formulas~\cite[Chapters I--II]{SaffTotik1997}.  They can also be checked directly:
with $d=c-r\cos\theta$, $\dd e_I=\dd\theta/\pi$, and the Cauchy transform
of $e_I$ is $((z-a)(z-2))^{-1/2}$; the second density has total mass one
and the difference of its logarithmic potential from $\log d$ is constant
on $I$.

Define the reference probability
\[
 \eta_{k,a}=(k+1)e_I-k\omega_{0,I}.
\tag{4.3}
\]
Thus
\[
 \dd\eta_{k,a}(d)=
 \left(k+1-\frac{k\sqrt{2a}}d\right)
 \frac{\dd d}{\pi\sqrt{(d-a)(2-d)}}.
\tag{4.4}
\]
The coefficient in parentheses is increasing, so $\eta_{k,a}\ge0$ if
and only if
\[
 a\ge2\left(\frac{k}{k+1}\right)^2.
\tag{4.5}
\]
Indeed, its minimum occurs at $d=a$ and is nonnegative precisely when
$k+1\ge k\sqrt{2/a}$, which is equivalent to (4.5).  It has total mass
$(k+1)-k=1$.  Let
\[
 D_0(a)=\frac{a+2+2\sqrt{2a}}4.
\]
On $I$, the equilibrium potential is $\log H$, while the balayage
potential is $\log d+\log H-\log D_0(a)$.  Hence
\[
 k\log d+\int_I\log|d-e|\dd\eta_{k,a}(e)=C(k,a),
\tag{4.6}
\]
where
\[
 C(k,a)=\log H+k\log D_0(a).
\tag{4.7}
\]
No sign assumption on $C(k,a)$ is made.

Let
\[
 W_0(x)=k\log|x|+\int_I\log|x-d|\dd\eta_{k,a}(d).
\]
Assume, as will be certified uniformly below, that the main component has
simple crossings
\[
 x_-<0<x_+<a,\qquad W_0'(x_-)<0<W_0'(x_+).
\tag{4.8}
\]
Set
\[
 \sigma_-=-\frac1{W_0'(x_-)},\quad
 \sigma_+=\frac1{W_0'(x_+)},\quad
 K_\pm=\sqrt{(a-x_\pm)(2-x_\pm)},
\tag{4.9}
\]
so $\sigma_->0$ and $\sigma_+>0$, and put
\[
 D=\frac{\sigma_-K_-}{a-x_-}+\frac{\sigma_+K_+}{a-x_+}.
\]
The adjoint measure is
\[
 d\xi(d)=
 \frac{D-\sigma_-K_-/(d-x_-)-\sigma_+K_+/(d-x_+)}
 {\pi\sqrt{(d-a)(2-d)}}\dd d.
\tag{4.10}
\]
If $N(d)$ denotes the numerator in (4.10), then the definition of $D$
gives $N(a)=0$, while
\[
 N'(d)=\frac{\sigma_-K_-}{(d-x_-)^2}
       +\frac{\sigma_+K_+}{(d-x_+)^2}>0.
\]
Thus $\xi$ is positive on every nonempty subinterval of $(a,2]$.
Moreover $N(d)=O(d-a)$ at $a$ and is bounded at $2$, so the square-root
singularities in (4.10) are integrable and $\xi$ is a finite measure.

\subsection{Convexity of the exact main width}

For a probability measure $\nu$ on $[1,2]$, let
\[
 T_\nu(u)=\inf\{d\in[1,2]:\nu([1,d])\ge u\},
 \qquad 0<u<1,
\]
be its nondecreasing quantile.  Write $T_0=T_{\eta_{k,a}}$ for the
reference quantile and $T$ for the quantile of the atomic target
$\sum_iq_i\delta_{d_i}$.  Thus $T$ equals $d_i$ on a consecutive
interval of length $q_i$.  For any such quantile $S$, let $M_k(S)$ be the
width of the component containing zero for
\[
 x\longmapsto k\log|x|+\int_0^1\log|x-S(u)|\dd u.
\]
For any nondecreasing quantile $S:[0,1]\to[1,2]$, define, for
$y<\operatorname*{ess\,inf}S$,
\[
 \Psi_S(y)=y\exp\left(\frac1k\int_0^1
        \log\left(1-\frac{y}{S(u)}\right)\dd u\right),\qquad
 R(S)=\exp\left(-\frac1k\int_0^1\log S(u)\dd u\right).
\tag{4.11}
\]
If $W_S(y)=k\log|y|+\int_0^1\log|y-S(u)|\dd u$, direct expansion gives
\[
 \frac1kW_S(y)=\log\frac{|\Psi_S(y)|}{R(S)}.
\tag{4.12}
\]

\begin{lemma}[Separation--contact criterion]\label{lem:separation}
Let $T$ be the quantile of an atomized target
$\sum_iq_i\delta_{d_i}$, with $1\le d_1<\cdots<d_N\le2$.  There is a
unique $y_c\in(0,d_1)$ at which $\Psi_T$ is critical, and
$R(T)\le\Psi_T(y_c)$.  Strict inequality means that the main component is
\emph{strictly separated} from the component at $d_1$; equality is the
\emph{contact} case, in which $y_c$ is their single common boundary point.
In both cases the main-component endpoints are $\Phi_T(-R(T))$ and
$\Phi_T(R(T))$, where $\Phi_T$ is the inverse branch of $\Psi_T$ at zero
and the latter value is the increasing real limit in the contact case.
\end{lemma}

\begin{proof}
For $0<y<d_1$, logarithmic differentiation of (4.11) gives
\[
 \frac{\Psi_T'(y)}{\Psi_T(y)}
 =\frac1y-\frac1k\sum_i\frac{q_i}{d_i-y}=:G(y).
\]
Now
\[
 G'(y)=-\frac1{y^2}-\frac1k\sum_i\frac{q_i}{(d_i-y)^2}<0,
\]
while $G(y)\to+\infty$ as $y\downarrow0$ and
$G(y)\to-\infty$ as $y\uparrow d_1$.  Hence $G$ has exactly one zero
$y_c$; the positive function $\Psi_T$ increases on $(0,y_c)$ and decreases
on $(y_c,d_1)$, tending to zero at both ends.

For $y<0$, the same logarithmic derivative is strictly negative, so
$\Psi_T'(y)>0$ because $\Psi_T(y)<0$.  Moreover $\Psi_T(y)\to-\infty$ as
$y\to-\infty$ and $\Psi_T(y)\to0$ as $y\uparrow0$.  Thus
$\Psi_T(y)=-R(T)$ has exactly one negative solution.

By atomization, the component containing zero and the component containing
$d_1$ are distinct.  Therefore some point of $(0,d_1)$ has $W_T\ge0$.
Equation (4.12) forces $R(T)\le\max_{(0,d_1)}\Psi_T=\Psi_T(y_c)$.
If the inequality is strict, $\Psi_T=R(T)$ has two positive solutions and
$W_T>0$ precisely between them; the first is the right endpoint of the main
component.  If equality holds, the solutions coalesce at $y_c$ and
$W_T<0$ on both adjacent punctured intervals, so $y_c$ is their common
boundary point.  Together with the unique negative solution, this proves
the assertions about the endpoints and the inverse branch.
\end{proof}

\begin{lemma}[Convex supporting inequality]\label{lem:convex-support}
Suppose that the reference $T_0$ is strictly separated and has the simple
crossings in (4.8).  If $T$ is an atomized separated or contact
target and $v=T-T_0$, then the directional derivative of the reference
width exists and
\[
 M_k(T)\ge M_k(T_0)+\dot M_k(T_0;v).
\tag{4.13}
\]
\end{lemma}

\begin{proof}
We prove the global supporting inequality, rather than assume a formal
first variation.  For a finite residual measure
$\sum_iq_i\delta_{d_i}$ put $\alpha_i=q_i/k$, so
$\sum_i\alpha_i=1/k$.  Formula (4.11) becomes
\[
 \Psi(y)=y\prod_i(1-y/d_i)^{\alpha_i},\qquad
 R=\prod_i d_i^{-\alpha_i}.
\]
For real $y<d_1$, all factors $1-y/d_i$ are positive, and
\[
 \frac1kW(y)=\log\frac{|\Psi(y)|}{R}.
\]
By Lemma~\ref{lem:separation}, the two main endpoints are $\Phi(-R)$ and
$\Phi(R)$, with the stated contact convention.  Lagrange inversion gives
\[
 \Phi(z)=\sum_{n\ge1}a_nz^n,\qquad
 a_n=\frac1n[y^{n-1}]\prod_i(1-y/d_i)^{-n\alpha_i}.
\tag{4.14}
\]
Consequently
\[
 M_k=\Phi(R)-\Phi(-R)=2\sum_{m\ge0}a_{2m+1}R^{2m+1}.
\tag{4.15}
\]
More explicitly,
\[
 a_nR^n=\frac1n
 \sum_{r_1+\cdots+r_N=n-1}
 \prod_i\frac{(n\alpha_i)_{r_i}}{r_i!}
 d_i^{-(n\alpha_i+r_i)}.
\tag{4.16}
\]
Here $(z)_0=1$ and
$(z)_m=z(z+1)\cdots(z+m-1)$ for $m\ge1$ is the rising Pochhammer symbol.
Every coefficient is nonnegative.  If
$m(d)=\prod_i d_i^{-\gamma_i}$ with $\gamma_i>0$, then for every vector
$v$,
\[
 \frac{v^T(\nabla^2m)v}{m}
 =\left(\sum_i\frac{\gamma_iv_i}{d_i}\right)^2
  +\sum_i\frac{\gamma_iv_i^2}{d_i^2}\ge0.
\]
Thus every finite odd partial sum in (4.15) is convex in the positive
coordinates.  Repeating coordinates according to multiplicity gives the
same statement on every common equal-mass quantile partition.

For completeness, we now pass from finite coordinate vectors to the
quantiles used below.  The masses $q_i$ of an empirical target are
rational.  Choose a common equal-mass partition which refines all of its
constant blocks, sample the reference quantile $T_0$ on the same cells,
and repeat each target value on the corresponding cells.  To see explicitly
that the changing dimension causes no problem, define
\[
 m_\ell(T)=\int_0^1T(u)^{-\ell}\dd u,
 \qquad R(T)=\exp\left(-\frac1k\int_0^1\log T(u)\dd u\right).
\]
The product in the coefficient formula (4.14) has the identity
\[
 \prod_i(1-y/d_i)^{-nq_i/k}
 =\exp\left(\frac nk\sum_{\ell\ge1}
 \frac{m_\ell(T)}\ell y^\ell\right).
\]
Thus, for fixed $n$, $a_n$ is a finite polynomial in
$m_1,\ldots,m_{n-1}$.  Under common-partition approximation these moments,
$R$, and their directional derivatives converge, because $T\ge1$ and
\[
 Dm_\ell(T)[v]=-\ell\int_0^1v(u)T(u)^{-\ell-1}\dd u,
 \qquad
 DR(T)[v]=-\frac{R(T)}k\int_0^1\frac{v(u)}{T(u)}\dd u.
\]
It follows that every fixed odd truncation and its directional derivative
converge as the mesh tends to zero.

Let $S_m$ denote the sum of the terms in (4.15) through degree $2m+1$.
Choose equal-mass step approximations $T_{0,N}$ to $T_0$, with $N$ a
multiple of the common denominator of the target masses.  Repeat each
target value on the same $N$ cells.  Finite-dimensional convexity gives
\[
 S_m(T)\ge
 S_m(T_{0,N})+DS_m(T_{0,N})[T-T_{0,N}].
\]
The displayed moment and derivative formulas are dominated uniformly by
$T_{0,N}\ge1$ and $|T-T_{0,N}|\le1$.  Hence, for fixed $m$, dominated
convergence as $N\to\infty$ yields
\[
 S_m(T)\ge S_m(T_0)+DS_m(T_0)[T-T_0].
\]
At a strictly separated reference, $R$ lies strictly inside the first
positive critical value of the inverse branch.  To justify that no nearer
complex singularity intervenes, recall Pringsheim's
theorem~\cite[Theorem IV.6]{FlajoletSedgewick2009}: a power series with
nonnegative coefficients and finite radius of convergence is singular at
the positive point equal to that radius.  The coefficients $a_n$ in
(4.14) are nonnegative, while the positive real inverse branch is analytic
up to its first critical value.  Thus the reference level is strictly
inside the disk of convergence.  On a smaller closed disk, the inverse
series and its material derivative converge locally uniformly by the
analytic implicit-function theorem.  Hence the series for both
$S_m(T_0)$ and $DS_m(T_0)$ converge absolutely.  Letting $m\to\infty$
in the preceding finite-truncation inequality gives the result when the
target is strictly separated.  Indeed, the same Pringsheim argument applies
to the target: Lemma~\ref{lem:separation} puts $R(T)$ strictly below its
first positive critical value, so its inverse series converges at
$\pm R(T)$.  Thus its nonnegative odd series converges to the exact
inverse-branch width, while the reference value and derivative converge
absolutely.

For a target at contact and $0<\lambda<1$, define
the Abel width
\[
 M_{k,\lambda}(T)=2\sum_{m\ge0}a_{2m+1}(T)
 [\lambda R(T)]^{2m+1}.
\]
For each finite odd truncation, convexity follows because its monomials are
precisely those in (4.16), multiplied by positive constants
$\lambda^{2m+1}$.  Apply the common-partition argument to that truncation
and then let its degree tend to infinity.  Monotone convergence applies on
the target side; at the strictly separated reference, both the value and
directional derivative converge absolutely.  Therefore
\[
 M_{k,\lambda}(T)\ge M_{k,\lambda}(T_0)
 +DM_{k,\lambda}(T_0)[T-T_0].
\]
As $\lambda\uparrow1$, the target value increases to $M_k(T)$: for
$\lambda<1$ it equals
$\Phi(\lambda R)-\Phi(-\lambda R)$, and the two real inverse branches
converge monotonically to the two level-$R$ endpoints.  This is precisely
the monotone form of Abel's boundary theorem, and it also follows directly
from monotone convergence of the nonnegative odd series.  Strict
separation gives absolute convergence of the reference value and material
derivative, so their limits may be taken term by term.  Letting
$\lambda\uparrow1$ proves (4.13) in the contact case.
\end{proof}

\subsection{The endpoint-corrected adjoint}

The pairing with $\xi$ is not, in general, equal to the derivative in
(4.13).  We now prove the exact endpoint correction and the one-sided
inequality that is needed.

For $0\le s\le1$, set $T_s=T_0+s(T-T_0)$.  In the reference spatial
coordinate $d=T_0(u)$, define
\[
 v(d)=T(u)-T_0(u).
\]
This is unambiguous almost everywhere because $\eta_{k,a}$ has no atoms.
Define $g$ as the derivative, at $s=0$, of the weighted potential evaluated
at the moving material point:
\[
 g(T_0(u))=
 \left.\frac{\dd}{\dd s}\right|_{s=0}
 \left\{k\log T_s(u)+
 \int_0^1\log|T_s(u)-T_s(w)|\dd w\right\}.
\]
The principal values below are symmetric principal values in the reference
coordinate.

\begin{lemma}[Endpoint-corrected adjoint]\label{lem:adjoint}
Under the hypotheses of Lemma~\ref{lem:convex-support}, one has the exact
endpoint correction
\[
 \dot M_k(T_0;v)-\int_Ig\dd\xi
 =-a_\pi v(2)\left(\frac{\sigma_-}{K_-}
                    +\frac{\sigma_+}{K_+}\right).
\]
For an atomic target, $v(2)=d_N-2\le0$, and consequently
\[
 M_k(T)\ge M_k(T_0)+\int_Ig\dd\xi.
\]
\end{lemma}

\begin{proof}
Write $d=c-r\cos\theta$, $0\le\theta\le\pi$, so
$\dd e_I=\dd\theta/\pi$.
Let
\[
 A(d)=k+1-\frac{k\sqrt{2a}}d,\qquad a_\pi=A(2)>0.
\]
For $j\in\{-,+\}$ there is a unique $0<\rho_j<1$ with
\[
 x_j=c-\frac r2(\rho_j+\rho_j^{-1}).
\]
Then
\[
 K_j=\frac{r(1-\rho_j^2)}{2\rho_j},\qquad
 P_\rho(\theta)=\frac{1-\rho^2}{1-2\rho\cos\theta+\rho^2}
 =1+2\sum_{n\ge1}\rho^n\cos(n\theta),
\tag{4.17}
\]
and $K_j/(d-x_j)=P_{\rho_j}(\theta)$.  Thus (4.10) is
\[
 d\xi=B(\theta)\frac{d\theta}{\pi},\qquad
 B(\theta)=D-\sum_j\sigma_jP_{\rho_j}(\theta),
\tag{4.18}
\]
with $D=\sum_j\sigma_jP_{\rho_j}(0)$.  This also proves $B(0)=0$ and
$B(\theta)>0$ for $0<\theta\le\pi$.

First assume that the spatial material velocity is smooth and put
$F(d)=A(d)v(d)$.  Differentiating the preceding definition of $g$ under
the integral gives
\[
 g(d)=\frac{k v(d)}d+\pv\int_I
 \frac{v(d)-v(e)}{d-e}\dd\eta(e).
\]
Differentiating (4.6) with respect to its spatial variable gives
\[
 \frac kd+\pv\int_I\frac1{d-e}\dd\eta(e)=0.
\]
Multiplying this identity by $v(d)$ and subtracting it from the preceding
display cancels every term containing $v(d)$ and gives
\[
 g(d)=\pv\int_I\frac{v(e)}{e-d}\dd\eta(e)
     =\pv\int_I\frac{F(e)}{e-d}\dd e_I(e).
\tag{4.19}
\]
At the two exterior crossings, ordinary differentiation and implicit
differentiation give
\[
 \dot M_k=-\sum_j\sigma_j\int_I\frac{F(d)}{d-x_j}\dd e_I(d).
\tag{4.20}
\]

Expand $F(c-r\cos\theta)=f_0+2\sum_{n\ge1}f_n\cos(n\theta)$, first as a
finite cosine polynomial.  Here $U_n$ denotes the Chebyshev polynomial of
the second kind, characterized by
$U_n(\cos\theta)=\sin((n+1)\theta)/\sin\theta$.  The finite Hilbert
transform identity
\[
 \pv\frac1\pi\int_0^\pi
 \frac{\cos(n\phi)}{\cos\theta-\cos\phi}\dd\phi
 =-U_{n-1}(\cos\theta)
\tag{4.21}
\]
follows from the zero transform of the constant function, the case $n=1$,
and the cosine recurrence.  Moreover
\[
 \frac1\pi\int_0^\pi P_\rho(\theta)U_{n-1}(\cos\theta)\dd\theta
 =\begin{cases}
 \dfrac{1+\rho^2-2\rho^{n+1}}{1-\rho^2},&n\text{ odd},\\[2mm]
 \dfrac{2\rho(1-\rho^n)}{1-\rho^2},&n\text{ even}.
 \end{cases}
\tag{4.22}
\]
This is obtained by inserting
$U_{2m}=1+2\sum_{\ell=1}^m\cos(2\ell\theta)$ and
$U_{2m-1}=2\sum_{\ell=1}^m\cos((2\ell-1)\theta)$.
We spell out the remaining coefficient calculation.  Put
\[
 \varepsilon_n=\frac1\pi\int_0^\pi
 U_{n-1}(\cos\theta)\dd\theta
 =\begin{cases}1,&n\text{ odd},\\0,&n\text{ even},\end{cases}
\]
and let $S_n(\rho)$ denote the left side of (4.22).  Equations
(4.18), (4.19), and (4.21) give
\[
 \int_Ig\dd\xi=-\frac2r\sum_{n\ge1}f_n
 \left(D\varepsilon_n-\sum_j\sigma_jS_n(\rho_j)\right).
\]
Poisson orthogonality and (4.20) give, on the other hand,
\[
 \dot M_k=-\sum_j\frac{\sigma_j}{K_j}
 \left(f_0+2\sum_{n\ge1}f_n\rho_j^n\right).
\]
Set
\[
 \Gamma=\sum_j\frac{\sigma_j}{K_j}
 =\frac2r\sum_j\frac{\sigma_j\rho_j}{1-\rho_j^2}.
\]
The coefficient of $f_0$ in the difference is $-\Gamma$.  If $n$ is
even, its coefficient is
\begin{align*}
&-\frac4r\sum_j\frac{\sigma_j\rho_j^{n+1}}{1-\rho_j^2}
 -\frac4r\sum_j\frac{\sigma_j\rho_j(1-\rho_j^n)}{1-\rho_j^2}\\
&\hspace{35mm}=-\frac4r\sum_j
 \frac{\sigma_j\rho_j}{1-\rho_j^2}=-2\Gamma.
\end{align*}
If $n$ is odd, use
\[
 D=\sum_j\sigma_jP_{\rho_j}(0)
 =\sum_j\sigma_j\frac{1+2\rho_j+\rho_j^2}{1-\rho_j^2}
\]
and the odd line of (4.22).  The coefficient becomes
\[
 \frac{2D}{r}-\frac2r\sum_j\sigma_j
 \frac{1+\rho_j^2}{1-\rho_j^2}
 =\frac4r\sum_j\frac{\sigma_j\rho_j}{1-\rho_j^2}=2\Gamma.
\]
Finally $d(\pi)=2$, and therefore
\[
 F(2)=f_0+2\sum_{n\ge1}(-1)^nf_n.
\]
The three coefficient formulas are exactly those of $-\Gamma F(2)$.
Thus
\[
 \boxed{\dot M_k-\int_Ig\dd\xi
 =-F(2)\sum_j\frac{\sigma_j}{K_j}
 =-a_\pi v(2)\left(\frac{\sigma_-}{K_-}
                    +\frac{\sigma_+}{K_+}\right).}
\tag{4.23}
\]

For an atomic target, $v$ is piecewise $C^1$ in the reference coordinate,
with finitely many jumps.  Apply (4.23) to the Poisson--Abel regularization
$F_\lambda=f_0+2\sum\lambda^nf_n\cos(n\theta)$.  For fixed
$\lambda<1$, the identity first holds for finite cosine truncations and
then for $F_\lambda$ by absolute convergence.  As $\lambda\uparrow1$, the
exterior integrals converge by domination because $d-x_\pm$ is bounded
away from zero.  The endpoint term also converges:
$F_\lambda(\pi)\to F(\pi)$ because the even periodic extension of
$\theta\mapsto F(c-r\cos\theta)$ is continuous at $\pi$.
For the interior pairing, the
zero principal value of the constant numerator gives, at every continuity
point,
\[
 g(c-r\cos\theta)=\frac1{\pi r}\pv\int_0^\pi
 \frac{F(c-r\cos\phi)-F(c-r\cos\theta)}
      {\cos\theta-\cos\phi}\dd\phi.
\]
We now give a direct uniform-integrability estimate; no commutation of the
finite Hilbert transform with Poisson convolution is used.  Write
$\widetilde F(\theta)=F(c-r\cos\theta)$.  Since
$v(d)=d_i-d$ on each spatial target block, $\widetilde F$ is piecewise
real analytic with finitely many interior jumps.  If those jumps are
$\Delta_\ell$ at $\theta_\ell$, decompose
\[
 \widetilde F(\theta)=F_c(\theta)+
   \sum_\ell\Delta_\ell\1_{[\theta_\ell,\pi]}(\theta).
\]
The function $F_c$ is continuous and piecewise real analytic.  Its cosine
coefficients $c_n$ satisfy $|c_n|\le Cn^{-2}$: integrate once by parts,
and use that the piecewise analytic derivative $F_c'$ has bounded variation,
whose Fourier coefficients are $O(n^{-1})$.

For a single jump term, its $n$th cosine coefficient is
$-\Delta_\ell\sin(n\theta_\ell)/(\pi n)$.  Using
$U_{n-1}(\cos\theta)=\sin(n\theta)/\sin\theta$ and summing
$\sum_{n\ge1}\lambda^n\cos(nt)/n=-\log|1-\lambda e^{it}|$ gives its
regularized transform exactly as
\[
 \frac{\Delta_\ell}{\pi r\sin\theta}
 \left\{\log|1-\lambda e^{i(\theta+\theta_\ell)}|
       -\log|1-\lambda e^{i(\theta-\theta_\ell)}|\right\}.
\]
The first logarithm is uniformly bounded near $\theta_\ell$; the second
has the only singularity.  Uniformly in $0<\lambda<1$,
\[
 \int_{|t|<\delta}|\log|1-\lambda e^{it}||\dd t
 \le C\delta(1+|\log\delta|).
\]
To verify this estimate, use
$|1-\lambda e^{it}|\asymp\sqrt{(1-\lambda)^2+t^2}$ for $|t|\le1$ and
integrate separately over $|t|\le1-\lambda$ and
$1-\lambda<|t|<\delta$; the part where the logarithm is positive is
uniformly bounded.  At $0$ and $\pi$ the two logarithms in braces have the
same endpoint value; the mean-value theorem gives a difference
$O(\sin\theta)$ uniformly in $\lambda$, so division by $\sin\theta$
remains bounded.

For the continuous part, put
$\delta_\theta=\min(\theta,\pi-\theta)$.  Since
$|\sin(n\theta)|\le\min(n\delta_\theta,1)$ and
$\sin\theta\ge2\delta_\theta/\pi$, the coefficient bound gives, uniformly
in $\lambda$,
\[
 \left|\sum_{n\ge1}\lambda^nc_nU_{n-1}(\cos\theta)\right|
 \le C\bigl(1+|\log\delta_\theta|\bigr).
\]
The last two bounds are integrable, $B(\theta)=O(\theta^2)$ at zero, and
$B$ is bounded at $\pi$.  Away from the finitely many jumps the displayed
series converge pointwise as $\lambda\uparrow1$; the logarithmic estimates
give uniform integrability on the omitted neighborhoods.  Hence the
pairings converge in $L^1(B(\theta)\dd\theta)$, which passes (4.23) to the
target velocity.
Since the
top target quantile is $d_N\le2$,
\[
 v(2)=d_N-2\le0,\qquad
 \dot M_k\ge\int_Ig\dd\xi.
\tag{4.24}
\]
Combining (4.13) and (4.24) gives the corrected supporting inequality
\[
 M_k(T)\ge M_k(T_0)+\int_Ig\dd\xi.
\tag{4.25}
\]
\end{proof}

\subsection{The block inequality}

Let the target quantile equal $d_i$ on consecutive half-open intervals
$I_i\subset(0,1)$ of length $q_i$.  Let
\[
 F_0(d)=\eta_{k,a}([a,d])
\]
be the reference distribution function and define the pullback
\[
 \widehat\xi=(F_0)_\#\xi.
\]
Thus
\[
 \int_0^1\varphi(u)\dd\widehat\xi(u)
 =\int_I\varphi(F_0(d))\dd\xi(d)
\]
for every bounded Borel function $\varphi$.  Put
\[
 r_i=\widehat\xi(I_i),\qquad
 \E(I_i)=\int_{I_i}\int_{I_i}
 \log|T_0(u)-T_0(v)|\dd v\dd\widehat\xi(u).
\tag{4.26}
\]
These mixed logarithmic integrals are absolutely convergent.  Indeed, in
angular coordinates
\[
 |d(\theta)-d(\phi)|
 =2r\left|\sin\frac{\theta+\phi}{2}\right|
       \left|\sin\frac{\theta-\phi}{2}\right|,
\]
and $A(\phi)$ and $B(\theta)$ are bounded on $[0,\pi]$.  Each of the two
resulting logarithms is integrable on $[0,\pi]^2$.  Thus Fubini's theorem
may be used below, including when two blocks share an endpoint.

\begin{proposition}[Block reduction]\label{prop:block}
Assume the hypotheses of Lemma~\ref{lem:adjoint}.  Put
$M_0=M_k(T_0)$, $R_0=\xi(I)$, and
$C_{\rm eff}=C+(L-M_0)/R_0$; here $R_0>0$ by the strict positivity of
$\xi$ on $(a,2]$.  If every nonempty quantile interval
$I'\subset(0,1)$ satisfies
\[
 \E(I')\ge qr\log\frac{qr}{2}+C_{\rm eff}r,
 \qquad q=|I'|,\quad r=\widehat\xi(I'),
\]
then the atomized target satisfies $\J_k\ge L$.
\end{proposition}

\begin{proof}
For $u\in I_i$, every off-block target difference and reference
difference has the same sign.  The tangent inequality
\[
 \frac{y-x}{x}\ge\log y-\log x\qquad(x,y>0)
\tag{4.27}
\]
applies to their absolute values and to the external-field coordinates.
Indeed, writing $g(u)$ as shorthand for $g(T_0(u))$, the unreduced first
variation is
\[
 g(u)=k\frac{T(u)-T_0(u)}{T_0(u)}
 +\pv\int_0^1
 \frac{[T(u)-T(v)]-[T_0(u)-T_0(v)]}
      {T_0(u)-T_0(v)}\dd v.
\]
For $v\in I_j$, $j\ne i$, (4.27) gives
\[
 \frac{(d_i-d_j)-[T_0(u)-T_0(v)]}{T_0(u)-T_0(v)}
 \ge\log|d_i-d_j|-\log|T_0(u)-T_0(v)|.
\]
For $v\in I_i$ the target difference is zero, so the quotient equals
$-1$ almost everywhere.  Applying (4.27) also to $d_i/T_0(u)$ and
integrating the preceding bounds yields
\begin{align*}
g(u)\ge{}&k\log d_i-k\log T_0(u)
 +\sum_{j\ne i}q_j\log|d_i-d_j|\\
&-\int_{(0,1)\setminus I_i}\log|T_0(u)-T_0(v)|\dd v-q_i.
\end{align*}
Now (2.9) makes the first and third terms equal to
$-q_i\log R_i$, while the platform identity (4.6) supplies exactly the
remaining constant $-C$.  Therefore
\[
 g(u)\ge-q_i\log R_i-C-q_i
 +\int_{I_i}\log|T_0(u)-T_0(v)|\dd v.
\tag{4.28}
\]
Integrating, summing, and using (4.25) yields
\[
 \J_k\ge M_0+\sum_i\left\{
 r_i[-q_i\log R_i-C-q_i]+\E(I_i)+2R_i\right\},
\tag{4.29}
\]
where $M_0=M_k(T_0)$.  As a function of $R_i>0$, the expression in braces
has derivative $-q_ir_i/R_i+2$ and positive second derivative
$q_ir_i/R_i^2$.  Its unique minimum is therefore
\[
 \E(I_i)-q_ir_i\log\frac{q_ir_i}{2}-Cr_i,
\tag{4.30}
\]
attained at $R_i=q_ir_i/2$.  Here $r_i>0$, because $\xi$ is positive on
every nonempty subinterval away from the single endpoint $a$.  Put
\[
 R_0=\xi(I)=\widehat\xi((0,1)),\qquad
 C_{\rm eff}=C+\frac{L-M_0}{R_0}.
\tag{4.31}
\]
$R_0$ is strictly positive.  It follows that the interval inequalities
\[
 \boxed{\E(I')\ge qr\log\frac{qr}{2}+C_{\rm eff}r}
\tag{4.32}
\]
for every nonempty quantile interval $I'\subset(0,1)$ (where
$q=|I'|$ and $r=\widehat\xi(I')$) imply, after summing over the target
blocks and using $\sum_i r_i=R_0$,
\[
 \J_k\ge M_0+(C_{\rm eff}-C)R_0=L.
\tag{4.33}
\]
This is the conclusion of the proposition.
\end{proof}

\section{A circle rearrangement theorem for every quantile interval}
\label{sec:circle}

Use the angular coordinate $d=c-r\cos\theta$ and write
\[
 \dd\eta=A(\theta)\frac{\dd\theta}{\pi},\qquad
 \dd\xi=B(\theta)\frac{\dd\theta}{\pi}.
\]
The function $d(\theta)$ is increasing.  The density factor
$A(d)=k+1-k\sqrt{2a}/d$ is increasing in $d$, so $A(\theta)$ is
increasing.  Moreover, (4.18) writes $B$ as a positive constant minus
positive multiples of $P_\rho(\theta)$, and every $P_\rho$ decreases on
$[0,\pi]$; hence $B(\theta)$ is also increasing.  Put
$a_\pi=A(\pi)>0$, $b_\pi=B(\pi)>0$.  For an angular interval $J$ define
\[
 Q=\int_J\frac{A(\theta)}{a_\pi}\dd\theta,\qquad
 R=\int_J\frac{B(\theta)}{b_\pi}\dd\theta.
\tag{5.1}
\]
Its physical masses are $q=a_\pi Q/\pi$ and $r'=b_\pi R/\pi$.  We use
$\E(J)$ for the mixed energy obtained by restricting the inner
$\eta$-integration and the outer $\xi$-integration in (4.26) to the
corresponding reference quantile interval.

\begin{theorem}[Circle block inequality]\label{thm:circle-block}
For every nonempty angular interval $J$, put
\[
 x_n(X)=\frac{\sin(nX)}{nX},\qquad
 h(X)=\log(2\pi)-\frac32
 +2\int_0^1(1-t)\log\frac{\sin(Xt)}{Xt}\dd t.
\]
Then $h(X)\ge0$ for $0<X\le\pi$, and
\[
\begin{aligned}
\frac{\E(J)}{qr'}-\log\frac{qr'}2-\frac{C_{\rm eff}}q
\ge{}&\log\frac{2H}{a_\pi b_\pi}+h(Q)+h(R)\\
&+\sum_{n\ge1}\frac{(x_n(Q)-x_n(R))^2}{n}
-\frac{\pi C_{\rm eff}}{a_\pi Q}.
\end{aligned}
\]
Consequently, strict positivity of the right-hand side proves the interval
hypothesis in Proposition~\ref{prop:block}.
\end{theorem}

\begin{proof}

The chord identity is
\[
 \frac{|d(\theta)-d(\phi)|}{H}
 =|e^{i\theta}-e^{i\phi}|\,|e^{i\theta}-e^{-i\phi}|.
\tag{5.2}
\]
Let $K(u)=\log|e^{iu}-1|$, and let
$\T=\R/(2\pi\Z)$ carry ordinary, nonnormalized Lebesgue measure.  On the
representative interval $[-\pi,\pi]$, extended periodically, define
\[
 f(\theta)=\1_J(|\theta|)\frac{A(|\theta|)}{a_\pi},\qquad
 g(\theta)=\1_J(|\theta|)\frac{B(|\theta|)}{b_\pi}.
\]
Because $A$ and $B$ are increasing, $0\le f,g\le1$, and (5.1) gives
\[
 \int_\T f(\theta)\dd\theta=2Q,\qquad
 \int_\T g(\theta)\dd\theta=2R.
\]
Splitting each circle integral into its positive and negative halves, the
four sign choices give two copies of $K(\theta-\phi)$ and two copies of
$K(\theta+\phi)$.  Hence (5.2) gives
\[
 \frac1{QR}\int_J\int_J\frac{AB}{a_\pi b_\pi}
 \log\frac{|d(\theta)-d(\phi)|}{H}\dd\theta\dd\phi
 =\frac1{2QR}\int_\T\int_\T f(\theta)g(\phi)
 K(\theta-\phi)\dd\theta\dd\phi.
\tag{5.3}
\]

Set $\mathcal L(u)=\log2-K(u)=-\log\sin(|u|/2)$ for $|u|\le\pi$.
It is nonnegative, integrable, even, and decreasing in $|u|$.  The circle
convolution--rearrangement theorem of
Baernstein~\cite{Baernstein1989} (also a circle case of the spherical
theorem in~\cite{BaernsteinTaylor1976}), applied first to
$\min(\mathcal L,N)$ and then by monotone convergence, says that its
cross-energy increases after symmetric decreasing rearrangement.  We give
the remaining bathtub step.  If $h$ is symmetric decreasing and
$0\le p\le1$ has circle mass $2Q$, then
\[
 \int_\T(\1_{[-Q,Q]}-p)(h-h(Q))\dd\theta\ge0;
\]
inside the centered arc both factors are nonnegative, and outside both are
nonpositive.  Since $\int(\1_{[-Q,Q]}-p)=0$, this proves
$\int ph\le\int_{-Q}^Qh$.  Also, the convolution of two nonnegative
symmetric decreasing circle functions is symmetric decreasing: by the
layer-cake representation it is enough to convolve two centered-arc
indicators, and their overlap length decreases with the circular distance
from the origin.  Let $f^*$ and $g^*$ be the symmetric decreasing
rearrangements.  Baernstein's theorem first gives
\[
 \int_\T f(\mathcal L*g)\le
 \int_\T f^*(\mathcal L*g^*).
\]
Because $\mathcal L*g^*$ is symmetric decreasing, the displayed bathtub
inequality permits us to replace $f^*$ by $\1_{[-Q,Q]}$.  The convolution
$\mathcal L*\1_{[-Q,Q]}$ is again symmetric decreasing, so a second
application replaces $g^*$ by $\1_{[-R,R]}$.  The $\mathcal L$-energy is
therefore at most that of these two centered-arc indicators.  Since the
additive $\log2$
term depends only on the two masses, reversing the sign shows that the
$K$-energy is minimized by these arcs.

Define their normalized energy by
\[
 \cA(Q,R)=\frac1{4QR}\int_{-Q}^Q\int_{-R}^R
 K(\theta-\phi)\dd\phi\dd\theta
 =-\sum_{n\ge1}
 \frac{\sin(nQ)\sin(nR)}{n^3QR}.
\tag{5.4}
\]
The series follows by first integrating the absolutely convergent Abel
series $-\sum_{n\ge1}\lambda^n\cos(nu)/n$ and then letting
$\lambda\uparrow1$ by dominated convergence for the integrable logarithmic
kernel.  Therefore the left side of (5.3) is at
least $2\cA(Q,R)$.  If $x_n(Q)=\sin(nQ)/(nQ)$, then termwise subtraction
gives
\[
 2\cA(Q,R)-\cA(Q,Q)-\cA(R,R)
 =\sum_{n\ge1}\frac{(x_n(Q)-x_n(R))^2}{n}\ge0.
\tag{5.5}
\]
Furthermore
\[
 \cA(Q,Q)=2\int_0^1(1-x)\log(2\sin(Qx))\dd x.
\]
Thus
\[
 h(Q):=\cA(Q,Q)-\log(Q/\pi)
 =\log(2\pi)-\frac32
 +2\int_0^1(1-x)\log\frac{\sin(Qx)}{Qx}\dd x.
\tag{5.6}
\]
The function $p(u)=\sin u-u\cos u$ satisfies $p(0)=0$ and
$p'(u)=u\sin u>0$ on $(0,\pi)$; hence
$(\sin u/u)'=-p(u)/u^2<0$.  It follows that $h$ decreases on
$(0,\pi]$.  Formula (5.4) gives $\cA(\pi,\pi)=0$, and therefore
$h(\pi)=0$ and $h(Q)\ge0$.

Returning to the physical masses in (5.1), equation (5.3) and the arc
comparison give
\[
 \frac{\E(J)}{qr'}\ge \log H+2\cA(Q,R).
\]
Also
\[
 \log\frac{qr'}2
 =\log\frac{a_\pi b_\pi QR}{2\pi^2},
 \qquad
 \frac{C_{\rm eff}}q=\frac{\pi C_{\rm eff}}{a_\pi Q}.
\]
Substituting (5.5) and
$\cA(Q,Q)=\log(Q/\pi)+h(Q)$, and then cancelling the factors
$QR/\pi^2$, gives the exact lower bound
\begin{align}
\frac{\E(J)}{qr'}-\log\frac{qr'}2-\frac{C_{\rm eff}}q
\ge{}&\log\frac{2H}{a_\pi b_\pi}+h(Q)+h(R)\\
&+\sum_{n\ge1}\frac{(x_n(Q)-x_n(R))^2}{n}
-\frac{\pi C_{\rm eff}}{a_\pi Q}.
\tag{5.7}
\end{align}
This formula includes all factors of $2$ and $\pi$.  It reduces (4.32) to
explicit scalar inequalities on a rectangle of normalized masses.
\end{proof}

\section{The terminal one-cut family and the constant
\texorpdfstring{$L$}{L}}\label{sec:onecut}

The zero-platform subfamily of (4.3) admits a convenient parameter $q$.
Set
\[
 s=\frac{1-q}{1+q},\qquad a^{\rm sh}=2s^2,\qquad
 H=\frac{2q}{(1+q)^2},\qquad
 D_0=\frac{2}{(1+q)^2}=\frac Hq.
\tag{6.1}
\]
Then $C(k,a^{\rm sh})=0$ precisely when
\[
 k=k(q)=-\frac{\log H(q)}{\log(H(q)/q)},\qquad
 A=\frac{k}{1+k}=\frac{\log H(q)}{\log q}.
\tag{6.2}
\]
Multiplying the residual probability (4.4) by $1-A$ and shifting back by
one gives the probability measure
\[
 \mu_q=A(q)\delta_{-1}+
 \frac{t+1-2A(q)s(q)}
 {\pi(t+1)\sqrt{(t-a(q))(1-t)}}\1_{(a(q),1)}\dd t,
\tag{6.3}
\]
where $a(q)=2s(q)^2-1$.  Its density is nonnegative if and only if
$A(q)\le s(q)$, because its numerator is increasing and its value at the
left edge is $2s(s-A)$.

Put $d(q)=\log(2/(1+q)^2)$ and $\tau=-\log q$.  Then
$A=1-d/\tau$, and $A\le s$ is equivalent to
\[
 f(q)=(1+q)d(q)+2q\log q\ge0.
\tag{6.4}
\]
Its derivative is
\[
 f'(q)=d(q)+2\log q
 =\log\frac{2q^2}{(1+q)^2}<0
\quad(0<q<3-2\sqrt2).
\tag{6.5}
\]
Since $f(0+)=\log2>0$ and the certified endpoint evaluation at
$3-2\sqrt2$ is negative, (6.4) has a unique zero $q_s$.

We now derive the exterior equation rather than merely recording it.  For
the shifted support $[2s^2,2]$ one has
\[
 c=1+s^2,\qquad r=2H,\qquad
 \rho_0=\frac{c-2s}{r}=q.
\]
For $u>1$ put
\[
 x^{\rm sh}=c-H(u+u^{-1}).
\]
Then $K_x=H(u-u^{-1})$, $\rho_x=1/u$, and
$c-x^{\rm sh}+K_x=2Hu$.  The Poisson expansion of the equilibrium and
balayage potentials therefore gives
\[
 W_0(x^{\rm sh})
 =k\log|x^{\rm sh}|+\log(Hu)-2k\log(1-q/u).
\]
Direct factorization gives
\[
 x^{\rm sh}=\frac{2}{(1+q)^2}
 \frac{(u-q)(1-qu)}u=D_0\frac{(u-q)(1-qu)}u.
\]
Using $\log H+k\log D_0=C=0$ and collecting logarithms, we obtain
\[
 W_0(x^{\rm sh})
 =-k\log\frac{u-q}{|1-qu|}+(k+1)\log u.
\]
Since $A=k/(k+1)$, the exterior zero equation is exactly
\[
 F_q(u):=A(q)\log\frac{u-q}{|1-qu|}-\log u=0.
\tag{6.6}
\]
On $1<u<q^{-1}$,
\[
 F_q'(u)=\frac{qu^2-B(q)u+q}{u(u-q)(1-qu)},
 \quad B(q)=1+q^2-A(q)(1-q^2).
\tag{6.7}
\]
For $q<q_s$, $A<s$ and hence $B>2q$.  The numerator has two reciprocal
positive roots.  More precisely, its value at $u=1$ is $2q-B<0$, whereas
at $u=q^{-1}$ it is
\[
 \frac{1+q^2-B}{q}=\frac{A(1-q^2)}q>0.
\]
Since the product of the two roots is one, exactly one critical point lies
in $(1,q^{-1})$.  The denominator in (6.7) is positive there, so
$F_q$ decreases before that point and increases afterward.  Because
$F_q(1)=0$ and $F_q(u)\to+\infty$ as $u\uparrow q^{-1}$, there is exactly
one nontrivial root $u_+\in(1,q^{-1})$.

At $q=q_s$ one has $A=s$ and $B=2q$, so the numerator in (6.7) is
$q(u-1)^2$.  Thus the nontrivial root merges with $u=1$, which justifies
the continuous definition $u_+(q_s)=1$.  On $u>q^{-1}$,
\[
 F_q'(u)=-\frac{A(q)(1-q^2)}{(u-q)(qu-1)}-\frac1u<0,
\]
while the endpoint limits are $+\infty$ and $-\infty$.  Thus there is
exactly one root $u_->q^{-1}$.

The shifted spatial coordinate used above is
$x^{\rm sh}=1+s^2-H(u+u^{-1})$.  After subtracting one, the root
$u_-$ gives the left endpoint and $u_+$ the right endpoint in (1.7).
Subtracting those two endpoint formulas gives exactly $\Lambda(q)$ in
(1.3).

\subsection{Certified global minimization}

We state exactly what is certified, then describe the exhaustive charts.

\begin{certificate}[One-cut global minimization]\label{comp:onecut}
On $0<q\le q_s$, $\Lambda$ has exactly one stationary point $q_*$,
$\Lambda'<0$ before it and $\Lambda'>0$ after it.  The enclosures in
Theorem~\ref{thm:main} hold.  In addition,
\[
 0.123630684649383<q_s<0.123630684649384.
\tag{6.8}
\]
\end{certificate}

\begin{proof}
The certificate uses only interval evaluations of the following analytic
equations and their differentiated forms.  Near $q=0$, set
\[
 \varrho=(-\log q)^{-1},\qquad z_\pm=qu_\pm.
\]
Then
\[
 G_\pm(\varrho,z)=A\log\frac{z-q^2}{|1-z|}-\log z-d(q)=0,
\tag{6.9}
\]
with $0<z_+<1<z_-$.  The length is
\[
 \Lambda=\frac{2(z_--z_+)}{(1+q)^2}
 +\frac{2q^2}{(1+q)^2}\left(\frac1{z_-}-\frac1{z_+}\right).
\tag{6.10}
\]
At $\varrho=0$, $(z_+,z_-)=(1/2,3/2)$.  On the complete tail
$0\le\varrho\le0.02$, direct interval substitution into (6.9), the strict
$G_z$ signs, and implicit differentiation give
\[
 -0.788<\frac{d\Lambda}{d\varrho}<-0.737.
\tag{6.11}
\]
To verify the stated continuation at $\varrho=0$, note that
$q=e^{-1/\varrho}\to0$, $A=1-\varrho d(q)\to1$, and (6.9) tends to
$-\log(2|1-z|)=0$.  Its two solutions are exactly $z=1/2$ and $z=3/2$.

At the soft edge, put $x=\log u_+$ and $y=x^2$.  After division by the
identically present root $x=0$, the plus equation is
\[
 D(q,y)=A-1+2A\sum_{n\ge1}q^n
 \frac{\sinh(n\sqrt y)}{n\sqrt y}=0.
\tag{6.12}
\]
For an exact, nonsingular evaluation of this series, define
\[
\begin{aligned}
 C(y)&=\cosh\sqrt y=\sum_{m\ge0}\frac{y^m}{(2m)!},\\
 S(y)&=\frac{\sinh\sqrt y}{\sqrt y}
       =\sum_{m\ge0}\frac{y^m}{(2m+1)!},\\
 b(q,y)&=\frac{qS(y)}{1-qC(y)},\qquad
 w(q,y)=y\,b(q,y)^2,\\
 \operatorname{atanhc}(w)&=\frac{\operatorname{atanh}\sqrt w}{\sqrt w}
       =\sum_{m\ge0}\frac{w^m}{2m+1},
\end{aligned}
\tag{6.12a}
\]
with the continuous values
$C(0)=S(0)=\operatorname{atanhc}(0)=1$.  Since
\[
 \sqrt w=\frac{q\sinh x}{1-q\cosh x}
\]
and
\[
 2\operatorname{atanh}\frac{q\sinh x}{1-q\cosh x}
 =\log\frac{1-qe^{-x}}{1-qe^x},
\]
the series in (6.12) is exactly
$b(q,y)\operatorname{atanhc}(w(q,y))$.  Thus the expression
actually evaluated is
\[
 D(q,y)=A(q)-1+2A(q)b(q,y)\operatorname{atanhc}(w(q,y)).
\tag{6.12b}
\]
In particular
\[
 D(q,0)=A(q)\frac{1+q}{1-q}-1,
\]
which vanishes exactly when $A(q)=s(q)$, that is, at $q=q_s$.  The
certificate proves $D_y>0$ on the soft chart and $0\le w<0.03$.  It
evaluates the three positive series in (6.12a) through $m=32$ and appends
explicit one-sided remainder bounds, including bounds for their first
derivatives.

In the bulk, proposed root boxes are first certified by opposite endpoint
signs and a strict $G_z$ sign.  A parametric interval-Newton contraction
then encloses $z_\pm(\varrho)$, after which
\[
 z'=-G_\varrho/G_z,\qquad
 z''=-\frac{G_{\varrho\varrho}+2G_{\varrho z}z'+G_{zz}(z')^2}{G_z}
\tag{6.13}
\]
is substituted into (6.10).  No floating-point root is accepted without
these sign and Newton checks.

The next two terminating decimals are retained at full precision because
they are exact rational centers of tests with radii $10^{-30}$ and
$10^{-48}$, respectively; shortening either center without changing the
verifier would change the certificate.  Put
\[
\begin{aligned}
 q_c&=
 0.02571553686652745032257637166391965344,\\
 \varrho_c&=-\frac1{\log q_c},\\
 c_s&=
 0.1236306846493834978974060904264788695442437883724.
\end{aligned}
\]
The exhaustive cover is:
\begin{center}
\begin{tabular}{p{35mm}p{72mm}c}
\toprule
\text{chart}&\text{exact domain}&\text{proved sign}\\
\midrule
$q\to0$ tail&$0\le\varrho\le1/50$&$\Lambda_\varrho<0$\\
bulk left&$1/50\le\varrho\le\varrho_c-1/200$&$\Lambda_\varrho<0$\\
stationary tube&$\varrho_c-1/200\le\varrho\le\varrho_c+1/200$&
$\Lambda_{\varrho\varrho}>0$\\
bulk right&$\varrho_c+1/200\le\varrho\le1/\log10$&$\Lambda_\varrho>0$\\
regular soft chart&$1/10\le q\le c_s-10^{-5}$&$\Lambda_q>0$\\
soft endpoint chart&$c_s-10^{-5}\le q\le q_s$&$\Lambda_q>0$\\
\bottomrule
\end{tabular}
\end{center}
The bulk right endpoint $\varrho=1/\log10$ is exactly $q=1/10$, so the bulk and
soft charts meet with no gap.  Inside the strictly convex tube, the
certificate proves opposite signs for $\Lambda_\varrho$ at
$\varrho=\varrho_c-10^{-30}$ and $\varrho=\varrho_c+10^{-30}$.
Continuity gives a stationary
point between them, and strict convexity gives at most one stationary
point in the entire tube.  The signs on the adjacent charts therefore
show that it is the unique stationary point on $(0,q_s]$ and is the global
minimum.  Outward substitution into (6.10) gives tighter boxes for $q_*$
and $L$; rounding their endpoints outward gives the more readable
enclosures in Theorem~\ref{thm:main}.

Finally, (6.5) proves that the function in (6.4) is strictly decreasing.
Opposite endpoint signs on
$[c_s-10^{-48},c_s+10^{-48}]$ therefore isolate its unique zero $q_s$
and give (6.8).  The endpoint soft chart is evaluated on the slightly
larger interval through $c_s+10^{-48}$ and is then intersected with
$q\le q_s$; thus the isolation box does not create a coverage gap.  All
these checks are implemented by the one-cut section of
\texttt{numerical\_verifier.py}.
\end{proof}

\section{Uniform certified calibration for all remaining
\texorpdfstring{$k$}{k}}\label{sec:calibration}

We now certify (4.32).  The formulas below also specify every scalar that
enters the program.

For $x<a$ put
\[
 K_x=\sqrt{(a-x)(2-x)},\qquad
 \rho_x=\frac{r}{c-x+K_x},\qquad
 \rho_0=\frac{c-\sqrt{2a}}r.
\tag{7.1}
\]
The exterior reference potential and its derivative are
\begin{align}
W_0(x)&=k\log|x|+\log\frac{c-x+K_x}{2}
       -2k\log(1-\rho_0\rho_x),\tag{7.2}\\
W_0'(x)&=\frac{k}{x}-\frac1{K_x}
 +\frac{2k\rho_0\rho_x}{K_x(1-\rho_0\rho_x)}.
\tag{7.3}
\end{align}
These follow from the Poisson expansion of (4.4).  More explicitly, for
$x<a$ the two Poisson sums are
\[
\begin{aligned}
 \int_I\log|x-d|\dd e_I(d)
   &=\log\frac{c-x+K_x}{2},\\
 \int_I\log|x-d|\dd\omega_{0,I}(d)
   &=\log\frac{c-x+K_x}{2}
     +2\log(1-\rho_0\rho_x).
\end{aligned}
\]
Substitution in $\eta=(k+1)e_I-k\omega_{0,I}$ gives (7.2).
Furthermore $K_x'=-(c-x)/K_x$ and $\rho_x'=\rho_x/K_x$, which gives
(7.3) by direct differentiation.  At the crossings define
$\rho_\pm=\rho_{x_\pm}$.  Then
\begin{align}
a_\pi&=1+\frac{2k\rho_0}{1+\rho_0},\tag{7.4}\\
b_\pi&=\frac{4\sigma_-\rho_-}{1-\rho_-^2}
      +\frac{4\sigma_+\rho_+}{1-\rho_+^2},\tag{7.5}\\
R_0&=\frac{2\sigma_-\rho_-}{1-\rho_-}
      +\frac{2\sigma_+\rho_+}{1-\rho_+},\tag{7.6}\\
Q_{\max}&=\frac\pi{a_\pi},\qquad
R_{\max}=\frac{\pi R_0}{b_\pi}.
\tag{7.7}
\end{align}
To derive these identities, evaluate
$A(\pi)=k+1-k\sqrt{2a}/2$ and express the result through $\rho_0$.
Next use
\[
 P_\rho(0)-P_\rho(\pi)=\frac{4\rho}{1-\rho^2},
\qquad
 \frac1\pi\int_0^\pi P_\rho(\theta)\dd\theta=1
\]
in (4.18).  This gives (7.4)--(7.6).  Finally, (5.1) applied to the
full angular interval gives (7.7).

For the affine range, $C_{\rm eff}<0$.  Put
\[
 \mathcal B=\log\frac{2H}{a_\pi b_\pi}
              +h(Q_{\max})+h(R_{\max}),\qquad
 P=-\frac{\pi C_{\rm eff}}{a_\pi}>0.
\tag{7.8}
\]
Write $\operatorname{sinc}x=\sin(x)/x$ for $x\ne0$ and
$\operatorname{sinc}0=1$.  Keeping the $n=1$ square in (5.7), it is
sufficient to prove
\[
 \mathcal B+\frac P Q+(\operatorname{sinc}Q-\operatorname{sinc}R)^2>0
\tag{7.9}
\]
on $0<Q\le Q_{\max}$, $0<R\le R_{\max}$.  If $Q\le R_{\max}$, the
left side is at least $\mathcal B+P/R_{\max}$.  The sinc function is
strictly decreasing on $(0,\pi)$ because
$\sin x-x\cos x>0$ there, as proved in
Section~\ref{sec:circle}.  Thus, if
$R_{\max}\le Q\le Q_{\max}$, monotonicity gives the lower bound
\[
 \mathcal B+\frac P Q+
 (\operatorname{sinc}Q-\operatorname{sinc}R_{\max})^2.
\tag{7.10}
\]
The derivative certified negative on 32 $Q$-subintervals is
\[
 -\frac{P}{Q^2}
 +2\bigl(\operatorname{sinc}Q-\operatorname{sinc}R_{\max}\bigr)
       \operatorname{sinc}'Q.
\]
It is then enough to check (7.10) at $Q_{\max}$.  For $h$, the verifier
uses
\[
 \cA(Q,Q)=-\sum_{n=1}^{80}\frac{\sin^2(nQ)}{n^3Q^2}+\mathcal R_{80},
 \qquad -\frac1{2\cdot80^2Q^2}\le\mathcal R_{80}\le0.
\tag{7.11}
\]
On the constant-edge range, the square sum is discarded and it suffices to
check $\mathcal B+P/Q_{\max}>0$.

\begin{samepage}
\begin{certificate}[Complete parameter cover]\label{comp:cover}
For every $k\ge29/20$, the reference choices below satisfy all positivity,
crossing, adjoint, and interval inequalities required by
Proposition~\ref{prop:block}:
\begin{center}
\begin{tabular}{@{}p{51mm}p{27mm}p{30mm}c@{}}
\toprule
\text{range}&\text{reference edge}&\text{certificate}&\text{strict scalar margin}\\
\midrule
$[36/25,21/10]$&$1153/500-k/4$&264 $k$-slabs&$>0.036$\\
$[21/10,21/5]$&$9/5$&840 $k$-slabs&$>0.036$\\
$[k(41542/10^6),\,k(26631/10^6)]$
 &$C=0$&20 $\varrho$-slabs&$>0.046$\\
$[k(26631/10^6),\,\infty)$
 &$C=0$&tail plus 26 $\varrho$-slabs&$>0.0056$\\
\bottomrule
\end{tabular}
\end{center}
Here the margin is a rigorous lower bound for the normalized left side of
(4.32), namely the left side of (5.7).  In particular the four ranges
overlap and cover every $k\ge29/20$.
\end{certificate}
\end{samepage}

\begin{proof}
For the first two rows the choices are
\[
 a(k)=\frac{1153}{500}-\frac k4,\qquad
 a(k)=\frac95,
\tag{7.12}
\]
respectively.  Each rational $k$-slab is evaluated with outward Arb balls.
At both slab endpoints, point roots merely propose intervals; opposite
$W_0$ endpoint signs, the strict $W_x$ sign, and the interval sign of
$x'=-W_k/W_x$ prove that the root box contains the unique branch for the
whole slab.  The program then checks (4.5), $\xi>0$, $C_{\rm eff}<0$,
(7.7), and the scalar reductions (7.9)--(7.11).  An unresolved sign aborts.

For the last two rows use (6.1)--(6.2).  The map $k(q)$ is strictly
decreasing.  Indeed, with $d(q)=\log(2/(1+q)^2)$,
\[
 k(q)=-\frac{\log q}{d(q)}-1,\qquad
 k'(q)=-\frac{(1+q)d(q)+2q\log q}
 {q(1+q)d(q)^2}<0
\]
for $0<q<q_s$, by (6.4).  Also $d(q)\to\log2$ and
$-\log q\to\infty$ as $q\downarrow0$, so $k(q)\to\infty$.  The refined
certificate covers
\[
 26631/10^6\le q\le41542/10^6,
\]
and proves $k(41542/10^6)<21/5$.  It keeps both self-gaps $h(Q)$ and
$h(R)$ in (5.7).  The simple terminal certificate covers
$0<q\le26631/10^6$; it proves the denominator corresponding to
$2H/(a_\pi b_\pi)$ is strictly below one.  In the terminal family
$C=0$, while Certificate~\ref{comp:onecut} gives $M_0\ge L$.
Therefore $C_{\rm eff}=(L-M_0)/R_0\le0$, and discarding its favorable
term is legitimate.

The arithmetic uses directed endpoints throughout.  Decimal-to-ball and
ball-to-decimal conversions are explicitly enlarged; the affine range
starts at $36/25<29/20$, so a representational endpoint cannot create a
gap.  The two terminal ranges meet exactly at $k(26631/10^6)$, and the
inequality $k(41542/10^6)<21/5$ gives a strict overlap with the
constant-edge range.  The complete implementation is the
uniform-calibration section of \texttt{numerical\_verifier.py}.
\end{proof}
\section{Completion and strictness of the lower bound}\label{sec:lower}

\begin{theorem}[Uniform lower bound]\label{thm:lower}
Every $f\in\mathcal P$ satisfies $|E_f|>L$.
\end{theorem}

\begin{proof}
Apply Lemma~\ref{lem:atomize}.  Its new sublevel set is contained in the
old one, so it suffices to bound the atomized polynomial.  The normalization
of Section~\ref{sec:reduction} either has $A=1$, in which case the length
is $2>L$ by the enclosure in Theorem~\ref{thm:main}, or gives
(2.7)--(2.10) with $k\ge1$.

For $1\le k\le29/20$, Proposition~\ref{prop:lowk} gives
$\J_k>2>L$.  For $k\ge29/20$, choose the reference from
Certificate~\ref{comp:cover}.  Theorem~\ref{thm:circle-block} and its
certified scalar margin prove the hypothesis of
Proposition~\ref{prop:block} for each of the finitely many target quantile
blocks; that proposition gives $\J_k\ge L$.

The inequality is strict.  In every calibrated finite range, each block has
$q_i>0$, and $r_i>0$ because the adjoint density is positive except at its
single left endpoint.  The certified scalar lower endpoint is strictly
positive; after multiplication by $q_ir_i$, every nonempty block contributes
a strict surplus.  The terminal ranges have the same strict circle surplus,
even when the one-cut reference itself has width $L$.  Hence
$\J_k>L$.  Finally, the original polynomial has length at least that of its
atomization, proving $|E_f|>L$.
\end{proof}

\section{Sharp recovery and nonattainment}\label{sec:recovery}

We construct finite polynomials whose lengths tend to $L$.  Directly
discretizing the zero platform (1.6) would create small additional negative
components, so we first insert a positive buffer.

\begin{lemma}[Positive-platform approximation]\label{lem:positive-buffer}
There are probabilities $\mu_A$ on $[-1,1]$, indexed by
$A>A_*$ sufficiently close to $A_*$, such that their logarithmic potential
is a positive constant on $[a_*,1]$, their negative set is one interval,
and
\[
 |E_{\mu_A}|\longrightarrow L\qquad(A\downarrow A_*).
\]
\end{lemma}

\begin{proof}
Fix $s=s_*$ and $a=a_*=2s^2-1$.  For $A>A_*$ sufficiently close to
$A_*$ define
\[
 \mu_A=A\delta_{-1}+
 \rho_A(t)\1_{[a,1]}(t)\dd t,
\quad
 \rho_A(t)=\frac{t+1-2As}
 {\pi(t+1)\sqrt{(t-a)(1-t)}}.
\tag{9.1}
\]
Let $\dd\omega(t)=\dd t/[\pi\sqrt{(t-a)(1-t)}]$.  The elementary arcsine
integrals
\[
 \int \dd\omega=1,\qquad
 \int\frac{\dd\omega(t)}{t+1}
 =\frac1{\sqrt{2(1+a)}}=\frac1{2s}
\]
give $\int\rho_A=1-A$, so $\mu_A$ is a probability.  Its numerator has
minimum $2s(s-A)$.  Since $q_*<q_s$, $A_*<s$; hence the density is strictly
positive for $A>A_*$ close enough.  We choose the neighborhood small enough
that $A<s$ throughout.

Write $V_A=V_{\mu_A}$ and put
$S(z)=\sqrt{(z-a)(z-1)}$, with $S(z)\sim z$ at infinity.  Off the
support, the arcsine Stieltjes transform and the elementary partial
fraction identity are
\[
 \int_a^1\frac{\dd\omega(t)}{z-t}=\frac1{S(z)},\qquad
 \frac1{(t+1)(z-t)}=\frac1{z+1}
 \left(\frac1{t+1}+\frac1{z-t}\right).
\]
Since $\rho_A(t)\dd t=[1-2As/(t+1)]\dd\omega(t)$, these formulas give
\begin{align*}
V_A'(z)
&=\frac A{z+1}+\frac1{S(z)}
-\frac{2As}{z+1}\left(\frac1{2s}+\frac1{S(z)}\right)\\
&=\frac{1-2As/(z+1)}{S(z)}.
\end{align*}
Thus, off the support,
\[
 V_A'(z)=\frac{1-2As/(z+1)}{S(z)}.
\tag{9.2}
\]
We also compute the platform value.  The probability
\[
 \dd\omega_{-1}(t)=\frac{2s}{t+1}\dd\omega(t)
\]
is the balayage of $\delta_{-1}$ onto $[a,1]$.  Translating the interval
identities (4.2)--(4.6), or equivalently inserting their Poisson series,
gives for $x\in[a,1]$
\[
 \int\log|x-t|\dd\omega(t)=\log\frac{1-a}{4},\qquad
 \int\log|x-t|\dd\omega_{-1}(t)=\log(x+1)+\log q_*.
\]
Because $\mu_A=\omega+A(\delta_{-1}-\omega_{-1})$, subtraction gives the
constant platform value
\[
 C(A)=\log\frac{1-a}{4}-A\log q_*.
\tag{9.3}
\]
At $A=A_*$ this is zero by $A_*\log q_*=\log H(q_*)$ and
$(1-a)/4=H(q_*)$.  Since $\log q_*<0$, $C(A)>0$ for $A>A_*$.

We record the complete zero count.  For $x<a$, the chosen branch has
$S(x)<0$.  On $(-\infty,-1)$, the numerator $x+1-2As$ of (9.2) is
negative and $(x+1)S(x)>0$, so $V_A'<0$.  The potential decreases from
$+\infty$ to $-\infty$ and has one simple zero.  On $(-1,a)$ the
derivative changes from positive to negative only at $x_0=2As-1<a$
(because $A<s$).  The potential rises from $-\infty$, then decreases to
the positive value $C(A)$, so it has exactly one simple zero, on the rising
branch.  On $(1,\infty)$, (9.2) is positive and there is no further zero.
Thus $E_{\mu_A}$ is one interval and its zero set consists of two points.

At $A=A_*$ this is exactly the terminal potential with $q=q_*$.  The
coordinate calculation in Section~\ref{sec:onecut} identifies its two crossings with
$u_-(q_*)$ and $u_+(q_*)$, hence with the two endpoints in (1.7); their
distance is $L$.  Since $q_*<q_s$, the derivative signs proved after
(6.7) show that both crossings are simple.  As $A\downarrow A_*$,
(9.2)--(9.3) vary jointly in $C^1$ on neighborhoods of those crossings.
The implicit-function theorem therefore gives
\[
 |E_{\mu_A}|\longrightarrow L.
\tag{9.4}
\]
\end{proof}

We now pass from $\mu_A$ to polynomials.  We use two elementary convergence
lemmas.

\begin{lemma}[Weak convergence implies $L^1$ convergence of potentials]
If $\nu_n\Rightarrow\nu$ are probabilities on $[-1,1]$, then for every
bounded interval $B$,
\[
 \norm{V_{\nu_n}-V_\nu}_{L^1(B)}\longrightarrow0.
\tag{9.5}
\]
\end{lemma}

\begin{proof}
The map $t\mapsto\log|\,\cdot-t|$ is continuous from $[-1,1]$ to
$L^1(B)$: this is the $L^1$ continuity of translations of the locally
integrable function $\log|x|$.  By compactness it is uniformly continuous.
Choose a continuous partition of unity $\phi_j$ and sample points $t_j$
so that
\[
 \sup_t\norm{\log|\,\cdot-t|-\sum_j\phi_j(t)
             \log|\,\cdot-t_j|}_{L^1(B)}<\eps.
\]
Weak convergence gives $\int\phi_j\dd\nu_n\to\int\phi_j\dd\nu$.
After integrating the displayed finite-rank approximation against
$\nu_n$ and $\nu$, the limsup in (9.5) is at most $2\eps$.  Let
$\eps\downarrow0$.
\end{proof}

\begin{lemma}[Stability of sign sets]
If $u_n\to u$ in measure on a finite-measure set and $|\{u=0\}|=0$, then
$\1_{\{u_n<0\}}\to\1_{\{u<0\}}$ in measure and in $L^1$.
\end{lemma}

\begin{proof}
For every $\delta>0$, up to null representatives,
\[
 \{\1_{u_n<0}\ne\1_{u<0}\}
 \subset\{|u|\le\delta\}\cup\{|u_n-u|\ge\delta\}.
\]
First let $n\to\infty$, then $\delta\downarrow0$.  For indicator
functions, convergence in measure of the symmetric difference is exactly
$L^1$ convergence.
\end{proof}

Every probability on $[-1,1]$ is weakly approximated by equal empirical
measures; for example, use the quantiles at $(j-1/2)/n$.  Thus choose
\[
 \nu_n=\frac1n\sum_{j=1}^n\delta_{t_{j,n}}\Rightarrow\mu_A,
 \qquad f_n(x)=\prod_{j=1}^n(x-t_{j,n}).
\]
The first lemma gives $L^1(-2,2)$ convergence of the potentials, hence
convergence in measure.  Every measure under consideration is supported on
$[-1,1]$, so its negative set is contained in $(-2,2)$; no length is lost
by working on this fixed interval.  The zero set of the positive-buffer
potential has measure zero, so the second lemma gives
$|E_{f_n}|\to|E_{\mu_A}|$.  Choose $A_m\downarrow A_*$ and then one
empirical approximation for each $m$ with length error below $1/m$.
Together with (9.4), this produces polynomials with $|E_f|\to L$.
Theorem~\ref{thm:lower} says every finite polynomial has length strictly
larger than $L$, so the infimum is not attained.

\section{The supremum}\label{sec:supremum}

This section reconstructs Tao's expansive-quantile proof of the sharp upper
bound~\cite{Tao1038}.  We include all analytic steps so that the only
computer-assisted input is the three scalar signs in
Certificate~\ref{cert:upper-scalars}.

Put
\[
 U_\mu(x)=-V_\mu(x)=\int\log\frac1{|x-t|}\dd\mu(t).
\tag{10.1}
\]

\begin{lemma}[Duality]\label{lem:duality}
Let $\mu$ be supported on a compact interval $K$.  If a nonzero finite
positive compactly supported measure $\nu$ satisfies $U_\nu(t)<0$ for
every $t\in K$, then $\supp\nu$ is not contained in
$\{U_\mu\ge0\}$.
\end{lemma}

\begin{proof}
On the product of the two compact supports, the kernel
$\log(1/|x-t|)$ is bounded below.  Add a constant to make it nonnegative
and apply Tonelli; after subtracting the same finite constant from both
iterated integrals,
\[
 \int U_\mu\dd\nu=\int U_\nu\dd\mu.
\]
If $\supp\nu\subset\{U_\mu\ge0\}$, the left side is nonnegative,
whereas the right side is strictly negative.  Indeed, $U_\nu$ is finite
and bounded below on $K$ and is negative at every point; the measurable
sets $\{U_\nu\le-1/m\}$ exhaust $K$, so one of them has positive
$\mu$-measure.  The possible diagonal value $+\infty$ is harmless under
Tonelli and is excluded on the right by the hypothesis $U_\nu<0$ on $K$.
\end{proof}

Fix $f\in\mathcal P$ and set $\mu=\mu_f$.  Then
$\{U_\mu>0\}=E_f$, while
\[
 \{U_\mu\ge0\}=\{x:|f(x)|\le1\}
\]
is a nonempty compact set: it is closed, contains every root, and is
contained in $[-2,2]$.  Assume for contradiction that
$|\{U_\mu>0\}|>B$, where $B=2\sqrt2$.  Let
$[t_0-1,t_0+1]$ be a length-two interval containing the roots, and let
$M_-,M_+$ be the extreme points of $\{U_\mu\ge0\}$.  Its span is at least
the measure of its subset $\{U_\mu>0\}$ and is therefore greater than
$B$.  Reflection lets us
assume $t_0-M_->\sqrt2$.  Translate $M_-$ to zero.  Then
\[
 t_0>\sqrt2,\qquad M_+>B,\qquad t_0-1>0.
\tag{10.2}
\]
The point $0$ lies strictly to the left of the root interval because
$t_0>\sqrt2$, so $U_\mu$ is continuous near $0$.  Extremality gives
$U_\mu(0)\ge0$; a strict inequality would persist on a short interval to
the left, contradicting the definition of $M_-$.  Hence $U_\mu(0)=0$.
If $t_0-1>1$, every root is more than
one unit from zero, giving $U_\mu(0)<0$; hence
\[
 \sqrt2<t_0\le2.
\tag{10.3}
\]

Set $a=t_0-1$ and $S=\{U_\mu>0\}\cap[a,\infty)$.  The part of the
positive set to the left of $a$ has measure at most $a$, so
$|S|>B-a$.  For $y\ge a$ let $F(y)=|S\cap[a,y]|$ and, for
$0\le u\le B-a$, define the right quantile
\[
 \phi(a+u)=\inf\{y\ge a:F(y)>u\}.
\tag{10.4}
\]
The set in (10.4) is nonempty even at $u=B-a$, because $|S|>B-a$.
Since
\[
 F(y)=\int_a^y\1_S(x)\dd x,
\]
$F$ is continuous, nondecreasing, and $1$-Lipschitz.  Put
$y_u=\phi(a+u)$.  By the definition of the infimum, $F(y)\le u$ for
$a\le y<y_u$, whereas there are points $z_n\downarrow y_u$ with
$F(z_n)>u$.  Continuity, together with $F(a)=0$ when $y_u=a$, therefore
gives $F(y_u)=u$.  Moreover, for every $\eps>0$ one may choose
$z\in(y_u,y_u+\eps)$ with $F(z)>u$; then
\[
 |S\cap(y_u,z]|=F(z)-F(y_u)>0,
\]
so $y_u\in\overline S$.  If $0\le u<v\le B-a$, the inclusion
$\{F>v\}\subset\{F>u\}$ gives $y_u\le y_v$, and the $1$-Lipschitz
property gives
\[
 v-u=F(y_v)-F(y_u)\le y_v-y_u.
\]
Thus the generalized inverse is expansive:
\[
 \phi(s')-\phi(s)\ge s'-s\quad(a\le s\le s'\le B).
\tag{10.5}
\]
Moreover, $u=F(y_u)\le y_u-a$ shows
$\phi(a+u)=y_u\ge a+u$, so $\phi\ge\mathrm{id}$.  We have also proved
$\phi([a,B])\subset\overline S\subset
\{U_\mu\ge0\}$.  Extend $\phi$ by $\phi(s)=s$ for $s<a$ and by
$\phi(s)=\phi(B)+(s-B)$ for $s>B$.  The resulting monotone map remains
expansive across both junctions and tends to $\pm\infty$ at the two ends.

For $t\in[a,t_0+1]$, choose a generalized inverse $s_0$ so that
\[
 \phi(s_0-)\le t\le\phi(s_0+).
\]
Such an $s_0$ exists by monotonicity and the two infinite end limits.  It
cannot be below $a$, where $\phi$ is the identity, and
$\phi\ge\mathrm{id}$ implies $s_0\le t$; hence
$a\le s_0\le t\le t_0+1$.  If $s<s_0$, expansiveness and the left
inequality for $s_0$ give
$t-\phi(s)\ge\phi(s_0-)-\phi(s)\ge s_0-s$.  If $s>s_0$, the right
inequality gives
$\phi(s)-t\ge\phi(s)-\phi(s_0+)\ge s-s_0$.  Thus
\[
 |t-\phi(s)|\ge|s_0-s|.
\tag{10.6}
\]
Since $a>0$, the extension fixes $0$, and the argument above gives
$U_\mu(0)=0$.  Therefore, with $K_\mu=\{U_\mu\ge0\}$,
\[
 \phi(\{0\}\cup[a,B])
 \subset\{0\}\cup\overline S\subset K_\mu.
\]
The set $K_\mu$ is closed.  Hence every finite measure $\nu$ supported on
$\{0\}\cup[a,B]$ satisfies $\supp(\phi_\#\nu)\subset K_\mu$.
Consequently, if such a positive measure $\nu$ satisfies
$U_\nu<0$ on $[a,t_0+1]$, then
\[
 U_{\phi_\#\nu}(t)\le U_\nu(s_0)<0
\]
on the entire root interval.  This contradicts Lemma~\ref{lem:duality}.
The integrated kernel inequality is justified
at a jump of $\phi$ by first truncating the logarithmic kernel above and
then using monotone convergence.  Equivalently, the single point
$s=s_0$ has no $\nu$-mass, since $U_\nu(s_0)<0$ is finite.  It remains
to construct $\nu$.

Put
\[
 F_0(x)=x-x\log|x|,\qquad F_0(0)=0,\qquad
 F_0'(x)=\log(1/|x|).
\]
The following three choices cover (10.3).  Every terminating decimal below
denotes the exact rational number represented by that decimal.

\paragraph{Case 1: $\sqrt2<t_0\le1.7624$.}
Take $\nu=\delta_0+C\delta_B$.  Its potential is convex on the root
interval, because
\[
 U_\nu''(t)=\frac1{t^2}+\frac{C}{(B-t)^2}>0.
\]
A convex function lies below the chord joining its endpoint values, so it
is negative throughout if it is negative at both endpoints.
At $a=t_0-1$ this requires
\[
 C>\frac{\log(1/a)}{\log(B-a)},
\]
and at $t_0+1$ it requires
\[
 C<\frac{\log(t_0+1)}{\log(1/(B-t_0-1))}.
\]
All denominators are positive.  Such $C$ exists exactly when
\[
 \log\frac1{t_0-1}\log\frac1{B-t_0-1}
 <\log(t_0+1)\log(B-t_0+1).
\tag{10.7}
\]

\paragraph{Case 2: $1.7987\le t_0\le2$.}
Take
\[
 \nu=\delta_0+0.7233\,\1_{[2.268742,B]}(x)\dd x.
\]
The support lies in $\{0\}\cup[a,B]$.  Its potential is
\[
 -\log t+0.7233\{F_0(B-t)-F_0(2.268742-t)\},
\tag{10.8}
\]
which is strictly negative for every $t\ge0.7987$.

\paragraph{Case 3: $1.7624<t_0<1.7987$.}
Take
\[
 \nu=\delta_0+0.192829\,\1_{[1.63,B]}\dd x
 +0.224\,\1_{[1.919,B]}\dd x+0.155\delta_B.
\]
The endpoint atom $B$ lies strictly to the right of the root interval.
The required potential is
\begin{align}
-\log t&+0.192829\{F_0(B-t)-F_0(1.63-t)\}\notag\\
&+0.224\{F_0(B-t)-F_0(1.919-t)\}
-0.155\log|B-t|,
\tag{10.9}
\end{align}
and is strictly negative on $0.7624\le t\le2.7987$.

\begin{certificate}[The three upper scalar inequalities]\label{cert:upper-scalars}
Inequality (10.7) holds on $(\sqrt2,1.7624]$; (10.8) is negative on the
stated half-line; and (10.9) is negative on the stated compact interval.
\end{certificate}

\begin{proof}
Decimal coefficients are parsed from strings into outward intervals, and
all domain endpoints are protected by explicit overlaps.  Near
$t_0=\sqrt2$, the difference in (10.7) is even about $\sqrt2$, has value and first derivative
zero there, and its second derivative is less than $-0.059$.  Adaptive
interval boxes cover the remaining range, with every upper bound below
$-6\cdot10^{-8}$.  For (10.8), adaptive boxes cover the finite range
through the exact rational endpoint $383/100>B+1$; beyond it
every distance is at least one and the $-\log t$ term is strict.  Every
finite upper bound is below $-10^{-9}$.  For (10.9), all interval upper
bounds are below $-10^{-4}$.  The artificial case junctions are enlarged by
$10^{-12}$, so no decimal conversion can leave a gap.  The executable proof
is the upper-scalar section of \texttt{numerical\_verifier.py}.
\end{proof}

\begin{theorem}[Sharp upper bound]\label{thm:upper}
Every $f\in\mathcal P$ satisfies $|E_f|\le2\sqrt2$, and equality holds for
$f=(x^2-1)^m$, $m\ge1$.
\end{theorem}

\begin{proof}
The preceding contradiction applies to the arbitrary polynomial fixed
above and proves $|E_f|\le2\sqrt2$.  For $f=(x^2-1)^m$,
\[
 |f(x)|<1\iff |x^2-1|<1,
\]
whose measure is $2\sqrt2$.
\end{proof}

Finally, for every nonconstant polynomial, $f-1$ and $f+1$ have only
finitely many real zeros.  Thus replacing $|f|<1$ by $|f|\le1$ never
changes the measured length.

\begin{proof}[Proof of Theorem~\ref{thm:main}]
Certificate~\ref{comp:onecut} proves the uniqueness of $q_*$ and the stated
outward enclosures for $q_*$ and $L$.  Theorem~\ref{thm:lower} gives
$|E_f|>L$ for every finite polynomial, while the diagonal empirical
construction in Section~\ref{sec:recovery} gives a sequence with lengths
tending to $L$; hence the infimum is $L$ and is not attained.
Theorem~\ref{thm:upper} gives the upper bound $2\sqrt2$ and the displayed
family of equality cases.  These four conclusions are exactly the claims
of Theorem~\ref{thm:main}.
\end{proof}

\appendix
\section{Numerical Verification}\label{app:verification}

The companion numerical verifier is the single human-readable file
\texttt{numerical\_verifier.py}, distributed with this paper.  It certifies
the explicit finite numerical comparisons appearing in the proof; it does
not replace any of the analytic reductions proved in the preceding sections.
More precisely, it verifies:
\begin{itemize}
\item the exact rational inequalities in the elementary range;
\item the one-cut root branches, soft-edge continuation, unique stationary
      point, and the enclosures for $q_s,q_*$, and $L$ in
      Certificate~\ref{comp:onecut};
\item every positivity, crossing, derivative, Fourier-tail, and scalar
      inequality in the exact parameter cover of
      Certificate~\ref{comp:cover}; and
\item the three scalar inequalities in
      Certificate~\ref{cert:upper-scalars}.
\end{itemize}
It never treats a sampled grid as a proof.

The interval principle used is standard: an expression evaluated with
directed outward rounding encloses the exact value for every point of its
input box~\cite{MooreKearfottCloud2009}.  The affine and Fourier parts use
Arb midpoint-radius balls~\cite{Johansson2017}.  The one-cut and upper
scalar parts use directed interval operations for addition, multiplication,
division, square root, exponential, logarithm, sine, and cosine.  The
refined terminal trilogarithm is evaluated by Arb.  Every infinite Taylor
or Fourier series has an explicit signed remainder bound.

Ordinary floating-point roots and optimizers are used only to propose
candidate boxes.  A box is accepted only after interval endpoint signs,
the required derivative orientation, and, where applicable, an interval
Newton contraction have been proved.  Adjacent parameter boxes have shared
or overlapping exact endpoints, and decimal conversions are enlarged
outward.  Thus the proof-relevant output consists of rigorous enclosures and
signs, not unverified floating-point approximations.
\begin{thebibliography}{99}

\bibitem{EHP1958}
P. Erd\H{o}s, F. Herzog, and G. Piranian,
\emph{Metric properties of polynomials},
J. Analyse Math. \textbf{6} (1958), 125--148.

\bibitem{Bloom1038}
T. F. Bloom,
\emph{Erd\H{o}s Problem \#1038},
Erd\H{o}s Problems, accessed 13 July 2026.
\url{https://www.erdosproblems.com/1038}

\bibitem{Bloom1038Discussion}
Erd\H{o}s Problems contributors,
\emph{Discussion of Erd\H{o}s Problem \#1038},
accessed 13 July 2026.
\url{https://www.erdosproblems.com/forum/thread/1038}

\bibitem{Tao1038}
T. Tao,
\emph{Sublevel sets of logarithmic potentials},
unpublished note on the sharp upper bound in Erd\H{o}s Problem 1038,
20 December 2025.
\url{https://terrytao.wordpress.com/wp-content/uploads/2025/12/erdos-1038.pdf}

\bibitem{Baernstein1989}
A. Baernstein II,
\emph{Convolution and rearrangement on the circle}, Complex Variables
Theory Appl. 12 (1989), 33--37.
\url{https://doi.org/10.1080/17476938908814351}

\bibitem{BaernsteinTaylor1976}
A. Baernstein II and B. A. Taylor,
\emph{Spherical rearrangements, subharmonic functions, and
$*$-functions in $n$-space}, Duke Math. J. 43 (1976), 245--268.
\url{https://doi.org/10.1215/S0012-7094-76-04322-2}

\bibitem{FlajoletSedgewick2009}
P. Flajolet and R. Sedgewick,
\emph{Analytic Combinatorics}, Cambridge University Press, Cambridge,
2009, Theorem IV.6.

\bibitem{Johansson2017}
F. Johansson,
\emph{Arb: efficient arbitrary-precision midpoint-radius interval
arithmetic}, IEEE Trans. Comput. 66 (2017), 1281--1292.
\url{https://doi.org/10.1109/TC.2017.2690633}

\bibitem{MooreKearfottCloud2009}
R. E. Moore, R. B. Kearfott, and M. J. Cloud,
\emph{Introduction to Interval Analysis}, SIAM, Philadelphia, 2009.
\url{https://doi.org/10.1137/1.9780898717716}

\bibitem{SaffTotik1997}
E. B. Saff and V. Totik,
\emph{Logarithmic Potentials with External Fields}, Grundlehren der
mathematischen Wissenschaften 316, Springer, Berlin, 1997.
\url{https://doi.org/10.1007/978-3-662-03329-6}

\end{thebibliography}
\end{document}
